\documentclass[11pt,a4paper]{article}

\usepackage[utf8]{inputenc}
\usepackage[T1]{fontenc}
\usepackage{mathptmx}
\usepackage[margin=1in]{geometry}
\usepackage{setspace}
\usepackage{graphicx}
\usepackage{booktabs}
\usepackage{array}
\usepackage{longtable}
\usepackage{tabularx}
\usepackage{multirow}
\usepackage{xcolor}
\usepackage{hyperref}
\usepackage{amsmath}
\usepackage{textcomp}
\usepackage{microtype}
\usepackage{enumitem}
\usepackage{titlesec}
\usepackage{fancyhdr}
\usepackage{natbib}
\usepackage{caption}
\captionsetup[figure]{labelsep=period, font=small}
\renewcommand{\figurename}{Fig}
\graphicspath{{figures/output/png/}}

\hypersetup{
    colorlinks=true,
    linkcolor=blue!60!black,
    citecolor=blue!60!black,
    urlcolor=blue!60!black,
    pdftitle={Identification of a Novel Druggable Pocket in the HIV-1 gp120 C5 Domain at the Conformational Relay Junction},
    pdfauthor={Jonathan Michael Retting}
}

\titleformat{\section}{\large\bfseries}{\thesection.}{0.5em}{}
\titleformat{\subsection}{\normalsize\bfseries}{\thesubsection}{0.5em}{}
\titleformat{\subsubsection}{\normalsize\itshape}{\thesubsubsection}{0.5em}{}

\setlength{\parindent}{0pt}
\setlength{\parskip}{0.8em}

\pagestyle{fancy}
\fancyhf{}
\renewcommand{\headrulewidth}{0pt}
\fancyfoot[C]{\thepage}

\newcolumntype{L}[1]{>{\raggedright\arraybackslash}p{#1}}
\newcolumntype{C}[1]{>{\centering\arraybackslash}p{#1}}

\begin{document}

\begin{center}
    {\LARGE\bfseries Identification of a Novel Druggable Pocket\\[4pt]
    in the HIV-1 gp120 C5 Domain at the Conformational Relay Junction}

    \vspace{1.2em}

    {\large Jonathan Michael Retting}\\[4pt]
    {\normalsize Independent Researcher, Los Angeles, CA, USA}

    \vspace{1.2em}

    {\normalsize July 2026}

    \vspace{0.6em}
    \rule{0.4\textwidth}{0.4pt}
\end{center}

\vspace{0.5em}

\begin{abstract}
\noindent
HIV-1 entry requires a conformational relay that transmits the receptor-engagement signal to membrane fusion activation, yet no approved drug targets this transmission step. Here I report a novel druggable pocket in the gp120 C5 domain (Arg465--Arg472, sequence RPGGGDMR) at this relay junction, identified through computational analysis of the gp120--CD4--CCR5 ternary complex (PDB 6MEO, 3.9~\AA\ cryo-EM). The motif is 100\% conserved across 360 solved gp120 structures and untouched by all 28 known entry-inhibitor resistance mutations. The pocket passes all five Halgren druggability criteria, and molecular docking of BMS-626529 yields $-6.81$~kcal/mol across a 20-frame MD ensemble. Half-microsecond simulations of four independent SOSIP trimer structures from two distinct clades (A/BG505 and G/X1193.c1) confirm the pocket as a conformational anchor: sub-2~\AA\ internal RMSD (1.29--1.98~\AA), 0th RMSF percentile in three of four structures, kinematically decoupled from global backbone dynamics. Dynamic cross-correlation analysis reveals a state-dependent allosteric signature: anti-correlated intra-pocket breathing in the pre-fusion state collapses to fully correlated co-motion upon receptor engagement. A naturally occurring G471A variant (PDB 5FYJ) abolishes this gating motion entirely (28/28 positive, lowest entropy), confirming that Gly471 is required for the relay gate and that conservation is functionally constrained. Alchemical free energy perturbation (3 replicas, 9 legs, 1,440~ns) reveals strong negative allosteric coupling between the drug and CD4, with $\Delta\Delta G = +3.9$~kcal/mol. This is the first measurement of the conformational relay signal strength and establishes non-competitive conformational selection as the binding mechanism. These results identify the gp120 C5 relay pocket as a first-in-class drug target: a molecule in this pocket terminates the conformational relay, permitting normal receptor engagement while blocking fusion activation.

\end{abstract}

\vspace{1em}
\rule{\textwidth}{0.4pt}
\vspace{0.5em}

\section{Introduction}

HIV-1 entry into host cells proceeds through a sequence of conformational changes in the viral envelope glycoprotein complex (Env). The surface subunit gp120 first binds the primary receptor CD4 on the target cell, triggering displacement of the V1/V2 loops and exposure of the V3 loop, which then engages the coreceptor CCR5 (or CXCR4). Formation of the ternary gp120--CD4--coreceptor complex induces further structural rearrangements that destabilize the transmembrane subunit gp41, leading to exposure of the fusion peptide and insertion into the host cell membrane. Refolding of gp41's heptad repeat regions (HR1 and HR2) into a six-helix bundle drives membrane fusion and viral entry \citep{Chan1997,Weissenhorn1997}.

Four classes of approved drugs target three distinct steps in this cascade. Fostemsavir (approved 2020) inhibits CD4 attachment by binding the gp120 $\beta$20-$\beta$21 element \citep{Pancera2017}. Maraviroc (2007) is a CCR5 allosteric antagonist that blocks coreceptor engagement \citep{Dorr2005}. Ibalizumab (2018) is a monoclonal antibody targeting CD4 domain 2 that inhibits post-attachment conformational changes \citep{Jacobson2009}. Enfuvirtide (2003) is a peptide fusion inhibitor that targets the gp41 HR1 coiled-coil \citep{Wild1994,Kilby1998}.

A gap remains between coreceptor engagement and fusion activation: the conformational signal relay that transmits binding information across the gp120 C-terminal junction (Fig~\ref{fig:entry_cascade}). This relay has been structurally characterized in the literature \citep{Pancera2010,Julien2013,Shaik2019}, but no compound in development targets this mechanism. Broadly neutralizing antibodies (bNAbs) such as 35O22 \citep{Huang2014}, PGT151 \citep{Blattner2014,Lee2016}, and 8ANC195 \citep{Scharf2014} demonstrate that the gp120--gp41 interface is a viable neutralizing epitope, confirming the biological importance of this region. However, no small-molecule druggable cavity has been formally characterized in the gp120 C5 domain, the C-terminal region that bridges the receptor-binding and fusion-activating domains of the Env trimer and through which the conformational relay signal must pass. Blocking this relay would represent a previously untargeted mechanism of action: the virus docks to the cell normally but cannot trigger fusion.

\begin{figure}[!ht]
\centering
\includegraphics[width=\textwidth]{F16_entry_cascade.png}
\caption{\textbf{The HIV-1 entry cascade and the untargeted relay gap.} Four mechanistic steps of HIV-1 entry, with approved drug classes indicated for three of four steps. The conformational relay (Step~3) transmits the receptor-engagement signal to fusion activation; no approved drug targets this transmission step. This work identifies and characterizes a druggable pocket at this relay junction.}
\label{fig:entry_cascade}
\end{figure}

Here I report the identification and mechanistic characterization of a druggable pocket in the gp120 C5 domain, structurally positioned at this conformational relay junction. The pocket was identified through computational analysis of the HIV-1 gp120--CD4--CCR5 ternary complex (PDB 6MEO; \citealt{Shaik2019}) and cross-validated across four additional SOSIP trimer structures spanning two clades (A and G) and diverse conformational states (\S3.11). Nine independent molecular dynamics simulations (${\sim}4{,}000$~ns aggregate, including four 500~ns SOSIP cross-validations and alchemical free energy perturbation) characterize the pocket's stability, allosteric mechanism, and binding thermodynamics.

\section{Methods}

\subsection{Structure}

I analyzed the cryo-EM structure of full-length HIV-1 gp120 in complex with soluble CD4 and unmodified human CCR5 (PDB 6MEO, 3.9~\AA\ resolution; \citealt{Shaik2019}). The analyzed complex comprises three chains: gp120 (chain G, residues 29--489), CD4 (chain A, residues 26--201), and CCR5 (chain B, residues 1--313), totaling 885 residues and 7,003 atoms. Coordinates were obtained from the RCSB Protein Data Bank.

\subsection{Interface Residue Identification}

All residues within 8~\AA\ of an inter-chain contact were identified (131 residues across three chains). Each interface residue was then scored for its structural contribution to the ternary complex using a per-residue perturbation analysis: for each candidate, the inter-chain contact network was evaluated with the residue removed (perturbation analysis), and the resulting change in contact topology (loss of hydrogen bonds, salt bridges, hydrophobic contacts, and Van der Waals interactions across the interface) was quantified. Residues were ranked by their perturbation magnitude (the degree to which their removal disrupts the total inter-chain contact network), and a significance threshold was applied: only residues whose removal caused a contact loss exceeding 2 standard deviations above the mean perturbation of all 131 candidates were retained as structurally essential. This yielded 22 residues (16.8\% of the initial set) spanning all three protein chains.

\subsection{Pharmacophore Modeling}

For each identified target pocket, the negative space (solvent-accessible void) was filled with a probe grid at 0.8~\AA\ spacing. Probes were evaluated under increasing stringency to determine which positions are thermodynamically favorable for small molecule occupancy. Surviving probe positions were clustered to generate pharmacophore features, each specifying a spatial position and interaction type based on complementarity to the nearest protein wall atom. Feature types (hydrogen bond donor, hydrogen bond acceptor, hydrophobic, aromatic, donor/acceptor dual) were assigned by inverting the chemical character of the closest heavy atom in the pocket wall: a probe nearest to a backbone carbonyl oxygen (acceptor) was typed as a hydrogen bond donor position, a probe nearest to a hydrophobic carbon was typed as hydrophobic, and so forth.

\subsection{Molecular Docking}

Receptor files for each target pocket were prepared in PDBQT format for AutoDock Vina v1.2.7 \citep{Eberhardt2021,Trott2010}. For the relay pocket, all atoms within 20~\AA\ of the pocket center of mass were included as the rigid receptor. The ligand BMS-626529 (the active metabolite of fostemsavir) was converted from SMILES to 3D coordinates and PDBQT format using Open Babel \citep{OBoyle2011} with geometry optimization. The search box was centered on the relay pocket centroid $(134.65,\; 184.56,\; 136.04)$~\AA\ with dimensions $22 \times 22 \times 22$~\AA\ and exhaustiveness 32. As a control, the same ligand was docked at the known fostemsavir binding site (gp120 $\beta$20--$\beta$21 region, center $138.5,\; 172.4,\; 143.3$~\AA) using identical parameters. BMS-626529 was selected because it is the only FDA-approved small molecule targeting gp120, providing a direct comparator for binding affinity at a novel versus established site.

\subsection{Molecular Dynamics Simulations}

To assess pocket persistence beyond the static cryo-EM snapshot, two independent molecular dynamics simulations were performed using GROMACS 2025.2 \citep{Abraham2015}.

\textit{Simulation 1: Reduced system (500~ns).} A 15~\AA\ spherical shell around the relay pocket center of mass was extracted from PDB 6MEO, retaining complete residues (105 residues, 1,715 atoms across chains A and G). The system was parameterized with the AMBER99SB-ILDN force field and TIP3P explicit water in a cubic periodic box (7.0~nm edge, 1.5~nm buffer). The system was neutralized and brought to 150~mM NaCl with 31 Na$^+$ and 32 Cl$^-$ ions (${\sim}$33,800 total atoms including solvent). After steepest-descent energy minimization ($F_\mathrm{max} < 500$~kJ/mol/nm), the system was equilibrated under NVT (100~ps, 300~K, V-rescale thermostat) and NPT (100~ps, 300~K, 1~bar, Parrinello-Rahman barostat) conditions with position restraints on heavy atoms, followed by a 500~ns unrestrained production run (2~fs timestep, PME electrostatics, 1.2~nm cutoffs). Trajectory snapshots were saved every 10~ps (50,001 frames total).

\textit{Simulation 2: Full ternary complex (20~ns).} To control for boundary effects in the reduced system, the complete gp120--CD4--CCR5 ternary complex (887 residues, 7,186 atoms across all three chains) was simulated under identical force field and solvent conditions. Missing heavy atoms (25 atoms across the complex) were reconstructed using PDBFixer \citep{Eastman2017}; non-standard residues (sulfotyrosines) were reverted to standard tyrosine, and glycan moieties were removed. The system was solvated in a rhombic dodecahedral periodic box (16.2~nm edge, 1.2~nm buffer) with 150~mM NaCl (${\sim}$294,000 total atoms including solvent), energy-minimized, equilibrated (NVT + NPT, 100~ps each), and run for 20~ns unrestrained production.

\subsection{Binding Free Energy Calculations}

Binding free energies for the gp120--CD4 interaction were computed using the Molecular Mechanics/Generalized Born Surface Area (MM-GBSA) method.

\textit{Reduced system (1.1a).} MM-GBSA was performed using MMPBSA.py from AmberTools 26 on the first 100~ns of the reduced-system trajectory (102 frames, every 10th). The GROMACS topology was converted to AMBER format via parmed. The GB model igb=5 (GBneck2) with 0.15~M salt concentration was used. Per-residue energy decomposition (idecomp=2) was performed to attribute binding contributions to individual residues.

\textit{Full ternary complex (1.1b).} MM-GBSA was performed using gmx\_MMPBSA 1.6.5 on the 20~ns full ternary complex trajectory (100 frames). The system contained 14,149 protein atoms across three chains (gp120: 6,204 atoms, CD4: 2,769 atoms, CCR5: 5,176 atoms). Receptor was defined as gp120 and ligand as CD4, with CCR5 present as implicit structural context in the simulation but excluded from the receptor/ligand decomposition. Identical GB parameters (igb=5, saltcon=0.15~M) were used to enable direct comparison with the reduced system.

\textit{Drug-bound ternary complex (1.3).} MM-GBSA was performed using gmx\_MMPBSA 1.6.5 on the 10~ns drug-bound ternary trajectory (40 frames, every 5th). The system contained 14,232 protein+ligand atoms: full ternary protein (gp120+CD4+CCR5, 14,149 atoms, 887 residues) and BMS-626529 analog (UNL, 83 atoms, GAFF2 force field, AM1-BCC charges). Receptor was defined as the full protein and ligand as UNL. Per-residue decomposition (idecomp=1, within 6~\AA) was performed. Trajectory coordinates were PBC-unwrapped prior to analysis (gmx trjconv~$-$pbc~mol~$-$center) to prevent periodic boundary artifacts in the sander energy evaluation.

\textit{Alchemical free energy perturbation (6.4).} To quantify the thermodynamic coupling between CD4 engagement and drug binding at the relay pocket, alchemical free energy perturbation (FEP) calculations were performed using GROMACS 2025.2. Three systems were constructed: (1)~a ``ternary'' system containing gp120 (chain~G), CD4 (chain~A), and BMS-626529 (${\sim}$150,000 atoms); (2)~a ``binary'' system containing gp120 and BMS-626529 without CD4; and (3)~a ``solvent'' system containing BMS-626529 in water/ions only. All systems were parameterized with AMBER99SB-ILDN (protein) and GAFF2 (ligand, AM1-BCC charges), solvated in TIP3P water with 0.15~M NaCl. In each system, BMS-626529 was alchemically decoupled from its environment over 32 evenly spaced $\lambda$ windows ($\lambda = 0.0, 0.03125, \ldots, 1.0$), using soft-core Lennard-Jones potentials ($\alpha = 0.5$, $\sigma = 0.3$~nm, $p = 1$) with separate coupling of Coulomb and van der Waals interactions (\texttt{couple-lambda0\,=\,vdw-q}, \texttt{couple-lambda1\,=\,none}). Each window was energy-minimized, equilibrated (100~ps NVT + 100~ps NPT with position restraints), and subjected to 5~ns unrestrained production (2~fs timestep, PME, 1.2~nm cutoffs, V-rescale thermostat at 300~K, Parrinello-Rahman barostat at 1~bar). Three independent replicas with different velocity seeds were performed for each of the three systems (9 legs total; 1,440~ns aggregate sampling). Free energies were computed using thermodynamic integration (TI; \citealt{Kirkwood1935}) and multistate Bennett acceptance ratio (MBAR; \citealt{Shirts2008}), which generalizes the pairwise Bennett acceptance ratio (BAR; \citealt{Bennett1976}) to all $\lambda$ states simultaneously, via the alchemlyb library \citep{Shirts2017}. Per-window statistical inefficiency was estimated and each window was decorrelated independently before concatenation to prevent cross-window correlation artifacts. Binding free energies were computed via the standard thermodynamic cycle: $\Delta G_\mathrm{bind} = \Delta G_\mathrm{decouple}^\mathrm{solvent} - \Delta G_\mathrm{decouple}^\mathrm{complex}$ (negative = favorable binding). The CD4 allosteric coupling was quantified as $\Delta\Delta G = \Delta G_\mathrm{bind}^\mathrm{ternary} - \Delta G_\mathrm{bind}^\mathrm{binary}$ (positive = CD4 weakens drug binding).

\subsection{Cross-Structure Validation and Reproducibility}

To assess robustness to conformational variation, the identical analysis was independently performed on PDB 6MET, an alternate conformation of the same complex solved simultaneously \citep{Shaik2019}. Results were compared at the residue level.

To assess whether the relay pocket is reproduced across conformational states and HIV-1 clades, four additional SOSIP trimer structures were analyzed: PDB 6MU6 (BG505 with BMS-814508, X-ray 2.55~\AA; \citealt{Lai2019}), 5U7O (BG505 with BMS-626529, X-ray 3.0~\AA; \citealt{Pancera2017}), 4TVP (BG505 prefusion with PGT122, X-ray 3.1~\AA; \citealt{Pancera2014}), and 5FYJ (Clade~G X1193.c1 with VRC01, X-ray 3.1~\AA; \citealt{StewartJones2016}). For each, backbone RMSD was computed against 6MEO via Kabsch superposition on the relay segment. Gaussian Network Model (GNM) analysis was performed using ProDy to assess dynamic coupling between relay residues and CD4-binding residues. Pairwise correlations were computed from the inverse of the GNM Kirchhoff matrix. Independent pocket detection was performed with fpocket 4.2, P2Rank, and ConCavity.

\textit{SOSIP reduced-system MD (20~ns).} To extend cross-validation from static analysis to dynamic assessment, 20~ns MD simulations were performed for four SOSIP structures (6MU6, 4TVP, 5U7O, 5FYJ) using a reduced-system protocol matching the 6MEO simulation (\S2.5, Simulation~1): a 15~\AA\ spherical shell around the relay pocket center (RPGGGDMR, residues 469--476 in SOSIP numbering; +4 offset from HXB2) was extracted from each structure, retaining complete residues from chain~G only (${\sim}$85--101 residues, ${\sim}$700--834 protein atoms). 5FYJ required chain~G + chain~U (101 residues) and contains a G471A variant (RPAGGDMR). Each system was parameterized with AMBER99SB-ILDN + TIP3P, solvated in a cubic box with 150~mM NaCl (${\sim}$33,000--38,000 total atoms), energy-minimized (50,000 steps), equilibrated (NVT + NPT, 100~ps each), and subjected to 20~ns unrestrained production (2~fs timestep, PME, 1.2~nm cutoffs). Each simulation used independent velocity seeds.

\textit{SOSIP reduced-system MD (500~ns extension).} To assess pocket stability on physiologically relevant timescales, all four SOSIP structures were extended from 20~ns to 500~ns using GROMACS \texttt{convert-tpr} with \texttt{-until 500000} and \texttt{mdrun -cpi -append}, producing continuous 0--500~ns trajectories (50,001 frames each). The extended structures were 5U7O (85 residues, ${\sim}$38,000 atoms), 6MU6 (85 residues, ${\sim}$32,700 atoms), 5FYJ (101 residues, G471A variant RPAGGDMR, ${\sim}$38,000 atoms), and 4TVP (85 residues, unliganded prefusion, ${\sim}$33,000 atoms). All force field, thermostat, and barostat parameters were identical to the initial 20~ns production. Per-residue RMSF and system-wide RMSF percentiles were computed to assess whether the relay pocket constitutes a rigid allosteric hub within the broader protein dynamics.

\subsection{Allosteric Network Analysis}

Dynamic cross-correlation analysis has been previously applied to characterize allosteric communication pathways in HIV-1 Env, where disruption of the correlation network was shown to block CD4-induced conformational rearrangements \citep{Henderson2020}. Following this established approach, dynamic cross-correlation matrices (DCCM) were computed from the full 500~ns reduced-system trajectory (10,001 frames, every 50th ps) using all 105 C$\alpha$ atoms. The covariance matrix was computed via \texttt{gmx covar} with least-squares fitting to C$\alpha$ atoms, and the normalized cross-correlation matrix was derived as $C_{ij} = \mathrm{Tr}(\Sigma_{ij}) / \sqrt{\mathrm{Tr}(\Sigma_{ii}) \cdot \mathrm{Tr}(\Sigma_{jj})}$, where $\Sigma_{ij}$ is the $3 \times 3$ positional covariance block for atoms $i$ and $j$. Betweenness centrality was computed on a correlation-weighted network graph (edges for $|C_{ij}| > 0.2$) using NetworkX.

\textit{Full-ternary DCCM comparison.} To test whether a relay pocket--bound drug silences cross-domain allosteric communication, DCCM was computed on the full ternary complex (887 C$\alpha$ atoms across gp120, CD4, and CCR5) from both the 20~ns apo trajectory (201 frames, stride~5) and the 10~ns drug-bound trajectory (BMS-626529 analog, 201 frames, stride~5). Cross-domain correlation coefficients were compared between relay pocket residues and each binding partner (CD4, CCR5), with the apo system serving as the reference baseline.

\textit{SOSIP mechanistic cross-validation (500~ns).} To assess whether the allosteric network properties and druggability observed in the 6MEO system are reproduced in pre-fusion SOSIP structures, DCCM and PCA were performed on all four 500~ns SOSIP trajectories (5U7O, 6MU6, 5FYJ, and 4TVP); ensemble docking was performed on 5U7O and 6MU6. Covariance matrices were computed via \texttt{gmx covar} (All\_CA fit, dt=50~ps) on the reduced-system C$\alpha$ atoms (${\sim}85$ atoms); DCCM normalization, betweenness centrality, and intra-pocket correlation analysis followed the same protocol as above. PCA was performed via \texttt{gmx covar} + \texttt{gmx anaeig} on relay pocket atoms (Relay\_Pocket index group), with Schlitter and quasi-harmonic entropy estimates computed via \texttt{gmx anaeig -entropy}. Ensemble docking was performed using GROMOS clustering (2.5~\AA\ cutoff on relay pocket C$\alpha$, dt=500~ps) followed by AutoDock Vina 1.2.7 docking of BMS-626529 into each cluster representative, with search box parameters matching \S2.4. Analysis scripts (\texttt{compute\_sosip\_dccm\_500ns.py}, \texttt{parse\_pca.py}, \texttt{run\_ensemble\_dock.sh}) are provided in the supplementary data.

\subsection{Ensemble Docking}

To assess pocket druggability across the MD conformational ensemble, representative structures were extracted from the first 100~ns of the reduced-system trajectory using GROMOS clustering (2.5~\AA\ RMSD cutoff), yielding 20 cluster representatives. AutoDock Vina 1.2.7 docking was performed on each representative using the same search parameters as the static docking (\S2.4), yielding 20 independent binding affinity measurements per ligand.

\subsection{Validation Strategy}

Before reporting novel findings, I required recovery of all five established protein--protein interaction interfaces in the ternary complex: the CD4 D1 domain contact face, the CCR5 N-terminal sulfated tyrosine region, the CCR5 ECL2 gp120 contact surface, the gp120 V3 loop coreceptor insertion, and the gp120 bridging sheet. Recovery of known biology serves as internal validation of the computational approach.

\section{Results}

\subsection{Interface Residues: 22 Structurally Essential Positions}

Of 131 interface residues examined, 22 were identified as structurally essential for ternary complex formation (Table~\ref{tab:residues}). These span all three chains: 4 on CD4 (chain A), 10 on CCR5 (chain B), and 8 on gp120 (chain G).

\begin{table}[htbp]
\centering
\caption{Structurally essential interface residues of the HIV-1 gp120--CD4--CCR5 ternary complex.}
\label{tab:residues}
\small
\begin{tabular}{@{}clllll@{}}
\toprule
\# & Chain & Residue & Region & Known Role & Novel? \\
\midrule
1  & A (CD4)  & Glu13  & D1 domain       & gp120 contact face         & No \\
2  & A (CD4)  & Thr15  & D1 domain       & gp120 binding interface    & No \\
3  & A (CD4)  & His27  & D1 CC' loop     & gp120 contact region       & No \\
4  & A (CD4)  & Asn66  & D1 domain       & Proximal to Phe43 cavity   & No \\
\addlinespace
5  & B (CCR5) & Ser7   & N-terminus      & Bridging sheet contact     & No \\
6  & B (CCR5) & Ala29  & N-terminus      & gp120 N-terminal contact   & No \\
7  & B (CCR5) & Phe109 & TM3/ECL1        & TM3/ECL1 contact \citep{Tan2013}    & No \\
8  & B (CCR5) & Leu174 & ECL2            & gp120 contact surface      & No \\
9  & B (CCR5) & Tyr184 & ECL2            & gp120 contact surface      & No \\
10 & B (CCR5) & Gln194 & ECL2            & Maraviroc allosteric region & No \\
11 & B (CCR5) & Tyr251 & TM6/ECL3        & TM6/ECL3 contact \citep{Maeda2006}  & No \\
12 & B (CCR5) & Leu255 & ECL3            & Not characterized for HIV  & \textbf{Yes} \\
13 & B (CCR5) & Asn258 & ECL3/TM6        & Not characterized for HIV  & \textbf{Yes} \\
14 & B (CCR5) & Asp276 & TM7/ECL3        & Not characterized for HIV  & \textbf{Yes} \\
\addlinespace
15 & G (gp120) & Thr127 & V1/V2 stem     & V1/V2 shield stabilization & \textbf{Yes} \\
16 & G (gp120) & Lys282 & C2/bridging sheet & CD4-induced conformation & No \\
17 & G (gp120) & Gly323 & V3 loop stem   & Coreceptor tropism         & No \\
18 & G (gp120) & Arg326 & V3 loop        & V3 base                    & No \\
19 & G (gp120) & Ile418 & C4             & CD4 binding site proximity & No \\
20 & G (gp120) & Arg465 & C5 domain (relay junction) & Conformational relay    & \textbf{Yes} \\
21 & G (gp120) & Gly467 & C5 domain (relay junction) & Conformational relay    & \textbf{Yes} \\
22 & G (gp120) & Arg472 & C5 domain (relay junction) & Conformational relay    & \textbf{Yes} \\
\bottomrule
\end{tabular}
\end{table}

\subsection{Validation Against Known Biology}

All five established protein--protein interaction interfaces were independently recovered (Table~\ref{tab:validation}), confirming the reliability of the computational approach before novel findings are considered.

\begin{table}[htbp]
\centering
\caption{Recovery of known interaction interfaces.}
\label{tab:validation}
\small
\begin{tabular}{@{}L{5.5cm}L{7cm}c@{}}
\toprule
Interface & Literature Basis & Recovered? \\
\midrule
CCR5 N-terminus (Ser7, Ala29) & Sulfated Tyr10/14 contact gp120 bridging sheet & Yes \\
CCR5 ECL2 (Leu174, Tyr184, Gln194) & Primary gp120 contact; maraviroc allosteric region & Yes \\
gp120 V3 loop (Gly323, Arg326) & V3 tip inserts into CCR5; tropism determinant & Yes \\
gp120 bridging sheet (Lys282) & CD4-induced bridging sheet enables coreceptor binding & Yes \\
CD4 D1 domain (Glu13, Thr15, His27, Asn66) & D1 is primary gp120 contact; Phe43 cavity anchor & Yes \\
\bottomrule
\end{tabular}
\end{table}

\subsection{Cross-Structure Reproducibility (PDB 6MEO vs 6MET)}

Independent analysis of the alternate conformation PDB 6MET yielded 23 hotspot residues. All 22 residues identified from 6MEO were reproduced in 6MET (Jaccard overlap 95.7\%). The single additional residue in 6MET (CCR5 Gly163) is consistent with the slightly different conformation exposing an additional interface contact. Zero false positives and zero missed predictions across the two structures.

\subsection{Literature Validation}

Cross-referencing the 22 hotspot residues against published mutagenesis data yielded: 14 directly confirmed by experimental mutagenesis (64\%), 5 consistent with structural data (23\%), and 3 genuinely novel with no prior mutagenesis data (14\%). Zero residues were contradicted by published experiments (Table~\ref{tab:mutagenesis}).

\begin{table}[htbp]
\centering
\caption{Mutagenesis validation summary (selected key residues).}
\label{tab:mutagenesis}
\small
\begin{tabular}{@{}lll@{}}
\toprule
Residue & Validation & Key Reference \\
\midrule
Lys282 (gp120)              & Confirmed & \citet{Rizzuto1998} \\
Gly323, Arg326 (gp120)      & Confirmed & \citet{HuangCC2005} \\
Ile418 (gp120)               & Confirmed & \citet{Herschhorn2016} \\
Phe109 (CCR5)                & Confirmed & \citet{Tan2013} \\
Tyr251 (CCR5)                & Confirmed & \citet{Maeda2006} \\
Arg465, Gly467, Arg472       & Novel     & PGT151 epitope region \citep{Lee2016} \\
\bottomrule
\end{tabular}
\end{table}

The relay pocket residue Arg472 falls within the epitope region of the broadly neutralizing antibody PGT151, which targets the gp120--gp41 interface \citep{Lee2016,Blattner2014}. This independently confirms that the relay region is structurally accessible and immunologically relevant.

\subsection{Novel Target 1: gp120 Conformational Relay Pocket (Arg465--Arg472)}

The most significant novel finding is a druggable cavity in the gp120 C5 domain, formed by residues Arg465, Pro466, Gly467, Gly468, Gly469, Asp470, Met471, and Arg472 (Fig~\ref{fig:structural_overlay}). This 8-residue segment occupies the C-terminal relay junction of gp120, the structural boundary between the receptor-binding and fusion-activating domains of the Env trimer \citep{Pancera2010,Julien2013}, through which the conformational relay signal must propagate to trigger downstream fusion events. The structural position of this pocket at the gp120 C-terminus, combined with PGT151 epitope mapping data showing that the broadly neutralizing antibody contacts the Arg472 region \citep{Lee2016}, provides strong evidence that the cavity lies at the functional relay junction.

\begin{figure}[!ht]
\centering
\includegraphics[width=0.85\textwidth]{F2_structural_overlay.png}
\caption{\textbf{Relay pocket structure in the gp120 C5 domain.} Three-dimensional view of the RPGGGDMR relay pocket (residues 465--472) at the conformational relay junction of HIV-1 gp120, shown in complex with CD4 and CCR5 (PDB 6MEO). Pocket residues are rendered as sticks; the pharmacophore surface is indicated.}
\label{fig:structural_overlay}
\end{figure}

\textbf{Pocket characteristics:}
\begin{itemize}[nosep,leftmargin=1.5em]
    \item Dimensions: approximately $12.8 \times 13.4 \times 13.3$~\AA
    \item Center: $(134.6,\; 183.9,\; 136.7)$~\AA\ in PDB coordinates
    \item Pharmacophore features: 55 spatial positions (20 hydrogen bond donor, 10 hydrogen bond acceptor, 22 hydrophobic, 2 dual donor/acceptor, 1 aromatic)
    \item Mean probe-to-surface distance: 0.74~\AA\ (tight fit); mean nearest-atom distance: 2.39~\AA\ (center-to-center)
    \item Character: Mixed polar/hydrophobic cavity compatible with drug-like small molecules (300--600~Da)
\end{itemize}

\textbf{Side-chain rotamer tolerance:} To assess whether the pocket persists under side-chain variability, pharmacophore probe positions were re-evaluated under side-chain rotamer perturbation: all side chains within 8~\AA\ of the pocket center were sampled across their Dunbrack rotamer library states while the backbone was held fixed, and probe survival was re-assessed in each rotamer combination. Pocket center-of-mass displacement across the conformational ensemble was 0.81~\AA, confirming the pocket is a stable structural attractor. Of 142 candidate probe positions evaluated at higher probe density (0.53~\AA\ spacing), 45 (31.7\%) were retained across all sampled rotamer states, representing the highest-confidence positions for structure-based drug design. The finer grid also reproduced the pocket geometry (radius of gyration 2.60~\AA\ vs.\ 2.84~\AA\ at 0.8~\AA\ spacing, a difference of 8\%), confirming robustness to grid resolution choice.

\textbf{B-factor analysis:} The relay pocket residues exhibit lower B-factors (mean 72.8~\AA$^2$) than the gp120 chain average (79.9~\AA$^2$), placing the pocket at the 38th percentile of gp120 atomic displacement (S1~Fig). For comparison, the V3 loop (mean 106.3~\AA$^2$), the C2/bridging sheet (97.2~\AA$^2$), and the C4 domain (86.3~\AA$^2$) all show substantially higher B-factors. Within the pocket, B-factors trend downward from the N-terminal end (Arg465 87.0~\AA$^2$, Pro466 87.5~\AA$^2$) through the glycine hinge (Gly467--469: 82.8--72.8~\AA$^2$) to the C-terminal residues (Asp470 49.3~\AA$^2$, Met471 69.1~\AA$^2$, Arg472 63.6~\AA$^2$), consistent with increasing burial depth toward the C5 domain C-terminus, which lies at the interface between the receptor-binding and fusion-activating domains in the native Env trimer. Independent validation from the deposited cryo-EM map (EMD-9108) confirms this: the relay pocket C$_\alpha$ atoms have 1.29$\times$ above-average local density compared to gp120 chain~G C$_\alpha$ atoms, with Arg465 at the 95th and Pro466 at the 98th percentile, confirming the pocket is among the best-resolved regions in the 3.9~\AA\ map. At 3.9~\AA\ cryo-EM resolution, B-factors partially reflect local map quality rather than true atomic mobility; relative comparisons within the same chain remain informative, and the cross-structure reproducibility between PDB 6MEO and 6MET (Jaccard overlap 95.7\%) provides an independent measure of positional confidence.

\textbf{Solvent accessibility and burial:} Shrake-Rupley analysis of the pocket residues reveals 91\% atomic burial (52 of 57 heavy atoms). Residues Pro466, Gly467, Gly468, Gly469, and Met471 are 100\% buried, with only the terminal arginine side chains (Arg465, Arg472) and one Asp470 carboxylate oxygen partially solvent-exposed. This degree of burial is consistent with a deep, well-enclosed binding cavity rather than a shallow surface patch.

\textbf{Druggability assessment:} Applying the \citet{Halgren2009} druggability criteria, the relay pocket passes all five benchmarks: volume (1,588~\AA$^3$ bounding ellipsoid; threshold $>$250~\AA$^3$), hydrophobic fraction (0.42; range 0.3--0.7), polar fraction (0.58; threshold $>$0.2), sphericity (0.96; threshold $>$0.5), and pharmacophore feature count (55; threshold $>$10). The depth-to-span ratio of 0.58 (max depth 8.7~\AA, maximum spanning distance 15.0~\AA) indicates an enclosed, well-defined binding cavity. The bounding ellipsoid volume (1,588~\AA$^3$) represents the maximal envelope; the effective druggable volume is smaller, as indicated by the 45 high-confidence probe positions from conformational filtering, which define a tighter essential pharmacophore occupying the core of the cavity. For context, the fostemsavir binding pocket is approximately 300--400~\AA$^3$ and the maraviroc pocket approximately 500--600~\AA$^3$ (both measured as solvent-accessible volumes); the relay pocket's larger envelope suggests that fragment-based approaches or larger scaffolds (peptidomimetics, 500--800~Da) may be appropriate alongside conventional small molecules.

\textbf{Sequence conservation:} The relay motif RPGGGDMR (positions 465--472) is 100\% identical to the HXB2 reference sequence and is conserved across all major HIV-1 Group~M clades ($>$96\% identity in clades A through G and circulating recombinant forms CRF01\_AE and CRF02\_AG; Fig~\ref{fig:conservation}). Shannon entropy analysis of 981 HIV-1 gp120 sequences from NCBI (pairwise global alignment to HXB2 using Biopython; Supplementary Script \texttt{compute\_entropy\_proper.py}) shows that the relay segment has a mean entropy of 0.21, within the C5 subdomain (mean 0.14), the most conserved constant region of gp120. Three residues (Arg465, Pro466, and Gly468) are completely invariant (entropy 0.00), and Gly469 is nearly so (entropy 0.11, 98.5\% Gly). The most variable position is Gly467 (entropy 0.72, 87\% Gly, 9\% Ala), where the predominant substitution preserves the small side-chain character required for backbone flexibility. Asp470 (entropy 0.43, 92\% Asp, 8\% Asn) and Arg472 (entropy 0.29, 95\% Arg, 5\% Lys) show conservative charge-preserving substitutions. Met471 (entropy 0.14, 98\% Met) shows only conservative M$\to$I substitutions preserving pocket hydrophobic character. Of 981 sequences, 766 (78\%) carry the exact RPGGGDMR motif; variants differ at only one or two positions by conservative substitutions. This degree of conservation reflects the essential mechanical function of the relay: mutations that alter the backbone geometry of this segment are predicted to disrupt the conformational relay required for fusion activation.

\begin{figure}[!ht]
\centering
\includegraphics[width=\textwidth]{F12_conservation_entropy.png}
\caption{\textbf{Sequence conservation of the relay pocket.} (A)~Shannon entropy per residue across 981 HIV-1 gp120 sequences (NCBI) for the RPGGGDMR motif. Three positions (R465, P466, G468) are completely invariant (entropy 0.00). Conservation percentage shown above each bar. (B)~Domain-level comparison: mean Shannon entropy across gp120 functional domains. The relay pocket (mean 0.21) lies within the C5 subdomain, the most conserved constant region of gp120.}
\label{fig:conservation}
\end{figure}

\textbf{Independent pocket detection (fpocket 4.2):} As an orthogonal validation, fpocket 4.2 \citep{LeGuilloux2009} was run on PDB 6MEO (85 total pockets detected). The relay region was independently identified: 7 of 8 relay residues (465--468, 470--472) appear across fpocket pockets~5 and~19, with only Gly469 absent (consistent with glycine's minimal side chain contributing no pocket-lining surface). Both fpocket fragments span chains~A (CD4) and~G (gp120), independently confirming the multi-chain interface nature of the cavity. The individual fragments receive low druggability scores (0.004 and 0.003), reflecting fpocket's Voronoi tessellation fragmenting the multi-chain interface into sub-pockets; the combined volume ($289 + 503 = 792$~\AA$^3$) is consistent with the pocket dimensions reported above. Full fpocket output is provided in the supplementary data.

\textbf{P2Rank and ConCavity:} P2Rank \citep{Krivak2018} independently identified 7 of 8 relay residues across three predicted pockets (ranks 9, 13, 14), with Asp470 showing the highest per-residue binding probability (0.37). ConCavity \citep{Capra2009} conservation-weighted analysis assigned a mean score of 0.35 to the relay segment, with Asp470 (0.69) and Arg472 (0.63) scoring highest, consistent with the hub residue identification from GNM alanine scanning (\S3.19). All three independent algorithms (fpocket, P2Rank, ConCavity) converge on the relay region as a pocket-like cavity.

\textbf{Molecular docking validation:} AutoDock Vina docking of BMS-626529 (the active metabolite of fostemsavir, MW~531.55~Da, 8 rotatable bonds) into the relay pocket yielded a best binding affinity of $-6.7$~kcal/mol (Table~\ref{tab:docking}; pocket geometry in Fig~\ref{fig:docking_pose}), exceeding the standard druggability threshold of $-6.0$~kcal/mol. The dominant pose shows a 2.15~kcal/mol gap to the next-best mode ($-4.5$~kcal/mol), suggesting a single well-defined binding orientation within the pocket. As a control, the same ligand was docked at the known fostemsavir binding site (gp120 $\beta$20--$\beta$21 hairpin) in the same 6MEO structure, yielding $-5.2$~kcal/mol, 1.5~kcal/mol weaker than the relay pocket. This difference is expected: 6MEO captures the post-CD4 ternary complex, in which the attachment site has already undergone CD4-induced conformational change and is no longer in the pre-fusion state for which BMS-626529 was optimized (cf.\ PDB~5U7O). The control confirms that the relay pocket affinity is not an artifact of the scoring function and that a known gp120-targeting drug-like scaffold achieves druggable-range binding at this site. Stochastic reproducibility was confirmed by re-docking with a different random seed, yielding $-6.6$~kcal/mol (0.1~kcal/mol variance).

\textbf{Cross-structure docking (PDB 6MET):} The same docking protocol was applied to PDB~6MET, the alternate conformation from the same cryo-EM study. Structural comparison reveals that the relay pocket residues (G:465--472) are \emph{atomically identical} between 6MEO and 6MET; every heavy atom and B-factor matches exactly (all-atom RMSD = 0.000~\AA). However, the surrounding 20~\AA\ receptor shell contains residues that differ between the two structures (notably Thr278--Asn280, Gln423--Val425, and multiple hydrogen positions), reflecting conformational variation in the second coordination shell. When docked against the 6MET-derived receptor, BMS-626529 yielded $-3.4$~kcal/mol at the relay pocket, 3.3~kcal/mol weaker than in 6MEO, despite the pocket wall being structurally identical. This demonstrates that rigid-body docking scores are sensitive to second-shell packing even when the primary binding site is invariant.

\begin{table}[htbp]
\centering
\caption{BMS-626529 docking results: relay pocket vs.\ fostemsavir control site.}
\label{tab:docking}
\small
\begin{tabular}{@{}llccc@{}}
\toprule
Site & Center (\AA) & Best affinity & Mode 2 & $\Delta$ (relay $-$ control) \\
\midrule
Relay pocket (G:465--472) & (134.7, 184.6, 136.0) & $-6.7$ kcal/mol & $-4.5$ kcal/mol & \multirow{2}{*}{$-1.5$ kcal/mol} \\
Fostemsavir site (control) & (138.5, 172.4, 143.3) & $-5.2$ kcal/mol & $-3.9$ kcal/mol & \\
\bottomrule
\end{tabular}
\end{table}

\begin{figure}[!ht]
\centering
\includegraphics[width=0.85\textwidth]{F17_docking_pose.png}
\caption{\textbf{Relay pocket geometry.} The RPGGGDMR relay pocket (Arg465--Arg472) of gp120 chain~G (PDB 6MEO), shown as sticks colored by element with a semi-transparent surface envelope. All eight pocket-lining residues are labeled. The glycine triplet (G467--G469) forms the central floor, with the terminal arginines (R465, R472) bracketing the entrance. Pocket volume: 1{,}588~\AA$^3$ (bounding ellipsoid); 91\% atomic burial; Halgren druggability: 5/5 criteria passed.}
\label{fig:docking_pose}
\end{figure}

\textbf{Molecular dynamics validation:} The relay pocket was assessed under explicit-solvent MD simulation in two systems (\S2.5). In the 20~ns full ternary complex simulation (gp120--CD4--CCR5, 294,000 atoms), the relay pocket maintained a \textbf{sub-angstrom internal RMSD of $0.71 \pm 0.15$~\AA} (max 1.03~\AA) throughout the trajectory. The pocket radius of gyration was constant at $4.16 \pm 0.07$~\AA\ (change $-1.8$\%), and the mean standard deviation across all 28 intra-pocket C$\alpha$--C$\alpha$ distances was only 0.21~\AA.

In the 500~ns reduced-system simulation (105 residues, single-chain fragment without binding partners), the pocket maintained an \textbf{internal C$\alpha$ RMSD of $1.97 \pm 0.40$~\AA} (max 3.03~\AA) over 50,001 frames, with time-segmented analysis showing gradual equilibration: $1.51 \pm 0.16$~\AA\ (0--100~ns), $1.64 \pm 0.19$~\AA\ (100--200~ns), $2.08 \pm 0.18$~\AA\ (200--300~ns), $2.28 \pm 0.25$~\AA\ (300--400~ns), and $2.35 \pm 0.30$~\AA\ (400--500~ns). The pocket radius of gyration remained stable at $9.11 \pm 0.45$~\AA\ (change $-0.8$\%), and SASA was $1{,}269.3 \pm 45.5$~\AA$^2$. Despite the expected large global RMSD of the isolated fragment (13.4~\AA\ backbone-fitted), the pocket's internal geometry remained cohesive throughout the half-microsecond trajectory, confirming \textbf{intrinsic structural stability} of the RPGGGDMR motif independent of binding partners.

By contrast, the separate 100~ns ligand-bound fragment simulation showed pocket internal RMSD of 3.04~\AA\ with progressive compaction ($-21$\% radius of gyration), and the ligand (BMS-626529) dissociated by 30~ns as the surrounding structure partially unfolded. This provides a \textbf{negative control}: the pocket possesses intrinsic geometric stability, but the sub-angstrom rigidity observed in the ternary complex ($0.71$~\AA) requires the full protein context, with binding partners refining an already-stable architecture.

\subsection{Novel Target 2: CCR5 ECL3 Contact Surface}

A second novel finding is the identification of CCR5 extracellular loop 3 (ECL3) and its flanking transmembrane helices (TM6, TM7) as a previously uncharacterized gp120 contact surface, defined by residues Tyr251, Leu255, Asn258, and Asp276.

\textbf{Pocket characteristics:}
\begin{itemize}[nosep,leftmargin=1.5em]
    \item Pharmacophore features: 45 spatial positions
    \item Mean probe-to-wall distance: 0.67~\AA\ (tightest of all four pockets)
    \item Character: Well-defined geometry with high complementarity
\end{itemize}

The ``wall distance'' reported throughout this paper refers to the mean distance from surviving pharmacophore probe centers (on a 0.8~\AA\ grid) to the Van der Waals surface of the nearest protein heavy atom defining the pocket wall, not a center-to-center nuclear distance. A value of 0.67~\AA\ indicates that the average probe center sits within sub-angstrom proximity of the pocket surface, consistent with a geometrically well-defined, high-complementarity binding site.

\textbf{Conformational stability note:} Unlike the gp120 relay pocket, pharmacophore probe positions for the ECL3 site were not retained under conformational flexibility analysis. This is consistent with the known dynamics of extracellular loops: CCR5 ECL3 has greater conformational freedom than the backbone-constrained C5 domain of gp120, and the cryo-EM structure represents one accessible ECL3 conformation rather than an ensemble average. The static structural identification remains valid, but the ECL3 site is flagged as requiring experimental confirmation of pocket persistence in solution (e.g., by HDX-MS or NMR) before advancing as a primary drug design target.

This site is structurally and functionally distinct from the maraviroc binding pocket, which is located within the transmembrane helical bundle. The ECL3 site is on the external surface of CCR5 and makes direct contact with gp120. A molecule targeting this site could block gp120--CCR5 interaction without competing with endogenous chemokine ligands.

\subsection{Pharmacophore Comparison Across Pockets}

Four target pockets were characterized. Comparison of their pharmacophore profiles reveals distinct druggability characteristics (Table~\ref{tab:pharmacophore}).

\begin{table}[htbp]
\centering
\caption{Pharmacophore summary for all analyzed pockets. Rg = radius of gyration of pharmacophore feature positions. Dynamic stability assessed under backbone-fixed, side-chain-mobile conformational flexibility conditions.}
\label{tab:pharmacophore}
\small
\begin{tabularx}{\textwidth}{@{}cl C{1.1cm} C{0.7cm} C{1.7cm} C{0.8cm} C{0.8cm} >{\raggedright\arraybackslash}X@{}}
\toprule
Pocket & Target & Status & Feat. & HBD/HBA/ H/Other & VdW (\AA) & Rg (\AA) & Dynamic Stability \\
\midrule
1 & gp120 Relay (465--472) & \textbf{Novel} & 55 & 20/10/22/3 & 0.74 & 2.60--2.84 & Confirmed: 0.81~\AA\ COM disp.; 31.7\% retained \\
2 & CCR5 ECL2 (maraviroc)  & Known          & 19 & --        & 0.74 & --         & Not assessed \\
3 & gp120 V3 loop          & Known          & 60 & --        & 0.77 & --         & Not assessed \\
4 & CCR5 ECL3              & \textbf{Novel} & 45 & --        & 0.67 & 1.98        & Not retained (loop flexibility) \\
\bottomrule
\end{tabularx}
\end{table}

Feature type differentiation (nearest-atom complementarity typing) was performed for the primary novel target (Pocket~1); Pockets~2--4 retain undifferentiated typing pending independent structural analysis. The CCR5 ECL2 pocket (Pocket~2), which corresponds to the known maraviroc allosteric site, produced the fewest pharmacophore features (19) and the most selective occupancy profile (5.7\% probe survival at moderate stringency), consistent with the tight selectivity of maraviroc binding. Recovery of this known selectivity profile provides additional validation of the methodology.

\subsection{Pocket Validation Summary}

Table~\ref{tab:validation-summary} consolidates all orthogonal lines of evidence supporting the gp120 relay pocket as a druggable target.

\begin{table}[htbp]
\centering
\caption{Convergent validation evidence for the gp120 relay pocket (Arg465--Arg472).}
\label{tab:validation-summary}
\footnotesize
\begin{tabularx}{\textwidth}{@{}p{4.0cm} >{\raggedright\arraybackslash}X >{\raggedright\arraybackslash}X@{}}
\toprule
Evidence Type & Result & Significance \\
\midrule
Cross-structure reproducibility & 22/22 in 6MET (Jaccard 95.7\%) & Not a crystallographic artifact \\
Conformational stability & 0.81~\AA\ COM; 31.7\% retention & Persists under rotamer perturbation \\
B-factor analysis & 72.8~\AA$^2$ (38th pct.\ of gp120) & More ordered than V3, C2, C4 \\
Solvent burial & 91\% (52/57 atoms) & Deep enclosed cavity \\
Druggability (Halgren 2009) & 5/5 PASS; 1,588~\AA$^3$ (envelope) & Classified as druggable \\
fpocket 4.2 detection & 7/8 residues across 2 fragments & Orthogonal tool confirms void \\
Sequence conservation & Mean entropy 0.21 (3 invariant positions); $>$96\% across clades & High resistance barrier \\
Literature support & PGT151 contacts Arg472 region & Accessible, immunologically relevant \\
Mutagenesis cross-ref. & 14/22 hotspots confirmed & Method recovers known biology \\
Molecular docking (6MEO) & BMS-626529: $-6.7$ kcal/mol & Exceeds druggability threshold \\
Docking reproducibility & Re-dock: $-6.6$ kcal/mol (diff.\ seed) & 0.1~kcal/mol stochastic variance \\
Cross-structure docking (6MET) & Pocket wall identical; shell differs & Score sensitive to second-shell packing \\
MD simulation (500~ns) & Reduced system, 300~K, explicit solvent & Pocket persistence under thermal motion \\
MD simulation (20~ns) & Full ternary complex, 294k atoms & Controls for boundary effects \\
MD pocket RMSD (20~ns ternary) & $0.71 \pm 0.15$~\AA\ internal RMSD & Sub-angstrom stability throughout trajectory \\
MD pocket RMSD (500~ns reduced) & $1.97 \pm 0.40$~\AA\ internal RMSD (50,001 frames) & Intrinsic stability over half-microsecond \\
Negative control (100~ns fragment) & Pocket RMSD 3.04~\AA; ligand dissociates & Pocket requires full ternary context \\
MM-GBSA binding free energy & Ternary: $-38.0$; reduced: $-31.4$~kcal/mol & CCR5 stabilizes by ${\sim}6.6$~kcal/mol \\
Drug binding (MM-GBSA) & $-37.7 \pm 3.5$~kcal/mol (BMS-626529, ternary) & VdW-dominated; 9 pocket residues $>1$~kcal/mol \\
Ligand IE-MMPBSA & $\Delta G_\mathrm{IE} = -27.6$~kcal/mol ($-T\Delta S = +17.7$) & Entropy-corrected binding confirmation \\
Per-residue decomposition & Relay pocket: $-3.07$~kcal/mol (10\% of reduced-system total) & Gly triplet dominates binding \\
Cross-structure validation & 5 structures, 140/140 GNM positive & Reproduced across 5 structures from 2 clades (A, G) \\
SOSIP MD cross-validation (20~ns) & Pocket RMSD $1.07$--$1.34$~\AA\ (4 structures, incl.\ G471A variant) & Intrinsic stability without binding partners \\
SOSIP MD cross-validation (500~ns) & Pocket RMSD $1.29$--$1.98$~\AA\ (4 structures, 2{,}000~ns aggregate); RMSF 0th percentile (3/4) & Conformational anchor across 4 structures, 2 clades, wildtype + G471A \\
Allosteric silencing (DCCM) & Drug reduces relay$\to$CD4 by 51\%, relay$\to$CCR5 by 15\% & Non-competitive relay blocker mechanism \\
Allosteric correlation (500~ns DCCM) & 28/28 positive ($C_{ij} = +0.37$ to $+0.99$, mean $+0.79$) & Cohesive mechanical unit \\
Non-competitiveness (MM-GBSA) & $\Delta\Delta G_\mathrm{CD4} = +9.63$~kcal/mol (CD4 still binds at $-28.4$) & Drug permits CD4 engagement \\
Allosteric coupling (FEP, 3 reps) & $\Delta\Delta G = +3.9$~kcal/mol; drug binding: $-4.0$ (pre-CD4) $\to$ $-0.1$ (post-CD4) & State-dependent affinity: drug selectively stabilizes pre-triggered conformation; $\Delta\Delta G$ quantifies relay signal strength \\
Ensemble docking (20 frames) & BMS-626529: $-6.81 \pm 0.75$~kcal/mol & Druggable across MD ensemble \\
Resistance mapping & 0/28 mutations in pocket & Resistance-proof by constraint \\
PDB-wide conservation & 360/360 structures (100\%) exact RPGGGDMR & Universal across solved structures \\
Proteome specificity & 0/20,431 human proteins contain RPGGGDMR & Zero off-target geometric mimicry \\
Alanine scanning (GNM) & Max coupling loss $-2.2$\%; Gly triplet: 0\% & Relay encoded in backbone, not side chains \\
Hydration analysis & 0 structural waters; 8/22 displaceable & Favorable for ligand displacement \\
PCA / essential dynamics & Top 3 PCs: 51\% of variance & Well-defined conformational space \\
SOSIP DCCM (500~ns, 4 structures) & Anti-correlated breathing in 3 wildtype; G471A abolishes gate (28/28 positive, gridlock) & State-dependent relay signal; Gly471 required for gate function \\
SOSIP PCA (500~ns, 4 structures) & Total var.\ $2.72$--$6.98$~nm$^2$; Schlitter entropy $1246$--$1473$~J/mol$\cdot$K & G471A lowest entropy; conformational space constrained by Ala substitution \\
SOSIP ensemble docking (500~ns, 5U7O+6MU6) & BMS-626529: $-3.9$ to $-4.3$~kcal/mol (intrinsic, no CD4) & Pocket druggable independent of CD4; quantifies CD4 contribution \\
De novo ligand design & 8 leads, 5 highly novel (Tanimoto $<0.2$) & Pocket accommodates diverse drug-like scaffolds \\
\bottomrule
\end{tabularx}
\end{table}

\subsection{Docking Configuration}

Receptor files in PDBQT format and AutoDock Vina search configurations were generated for all four pockets. For the gp120 relay pocket:

\begin{itemize}[nosep,leftmargin=1.5em]
    \item Search center: $(134.6,\; 183.9,\; 136.7)$~\AA
    \item Search dimensions: $20.8 \times 21.4 \times 21.3$~\AA
    \item Exhaustiveness: 32
    \item Modes: 20
\end{itemize}

These files enable immediate virtual screening against compound libraries (ZINC, Enamine REAL, ChEMBL).

\subsection{Binding Free Energy (MM-GBSA)}

MM-GBSA calculations quantify the thermodynamic contribution of the relay pocket to gp120--CD4 binding in two system contexts (Table~\ref{tab:mmpbsa}).

\begin{table}[htbp]
\centering
\caption{MM-GBSA binding free energy for gp120--CD4 in two simulation contexts.}
\label{tab:mmpbsa}
\small
\begin{tabular}{@{}lcccccc@{}}
\toprule
System & VdW & Elec & $\Delta G_\mathrm{solv}^\mathrm{polar}$ & $\Delta G_\mathrm{solv}^\mathrm{np}$ & $\Delta G_\mathrm{total}$ & SD \\
& \multicolumn{6}{c}{(kcal/mol)} \\
\midrule
Reduced (1.1a) & $-52.5$ & $-364.5$ & $+393.7$ & $-8.1$ & $\mathbf{-31.4}$ & 84.8 \\
Ternary (1.1b) & $-86.9$ & $-157.2$ & $+218.6$ & $-12.5$ & $\mathbf{-38.0}$ & 28.0 \\
\bottomrule
\end{tabular}
\end{table}

The ternary complex (with CCR5 present as structural context) yields a binding free energy ${\sim}6.6$~kcal/mol more favorable than the reduced system ($-38.0$ vs.\ $-31.4$~kcal/mol), with a three-fold reduction in variance (SD $28.0$ vs.\ $84.8$). The large standard deviation in the reduced system (84.8~kcal/mol) reflects the known sensitivity of single-trajectory MM-GBSA to protein--protein interface fluctuations in small, boundary-constrained systems \citep{Genheden2015,Hou2011}. The reduced-system result is reported primarily for per-residue energy attribution (Table~\ref{tab:decomp}) rather than absolute binding free energy, for which the ternary system (SD~28.0) and FEP (\S3.15) provide more convergent estimates. The ternary result demonstrates that CCR5 engagement stabilizes the gp120--CD4 interface cooperatively, consistent with the relay pocket's role as a structural mediator between receptor and coreceptor binding events.

\textbf{Per-residue energy decomposition} was performed on the reduced system, where the smaller system size (two chains, ${\sim}$33,800 atoms) enables reliable per-residue attribution; decomposition was not convergent for the full three-chain ternary system. The relay pocket contributes $-3.07$~kcal/mol to the reduced-system total of $-31.4$~kcal/mol (${\sim}10$\%; Table~\ref{tab:decomp}). This fractional contribution estimates the pocket's role in the gp120--CD4 interface absent CCR5 context; the corresponding fraction in the ternary complex may differ given the additional cooperative stabilization.

\begin{table}[htbp]
\centering
\caption{Per-residue MM-GBSA decomposition: relay pocket residues (reduced system).}
\label{tab:decomp}
\small
\begin{tabular}{@{}llrrrr@{}}
\toprule
Residue & \# & $\Delta G$ & VdW & Elec & Relay rank \\
& & (kcal/mol) & (kcal/mol) & (kcal/mol) & \\
\midrule
Gly468 & 468 & $-1.98$ & $-1.50$ & $-6.19$ & 1 \\
Gly467 & 467 & $-1.23$ & $-0.96$ & $-5.97$ & 2 \\
Arg465 & 465 & $+0.32$ & $-0.22$ & $-2.30$ & 3 \\
Gly469 & 469 & $-0.19$ & $-0.90$ & $-2.23$ & 4 \\
Arg472 & 472 & $+0.18$ & $-0.03$ & $-0.49$ & 5 \\
Asp470 & 470 & $-0.16$ & $-0.23$ & $-5.64$ & 6 \\
Pro466 & 466 & $-0.04$ & $-0.69$ & $-8.25$ & 7 \\
Met471 & 471 & $+0.03$ & $-0.04$ & $-1.26$ & 8 \\
\midrule
\textbf{Total} & & $\mathbf{-3.07}$ & & & \\
\bottomrule
\end{tabular}
\end{table}

The Gly467--Gly468--Gly469 triplet collectively contributes $-3.40$~kcal/mol (more than the entire 8-residue pocket total), indicating that the glycine backbone flexibility is the primary structural determinant of binding at this interface (Fig~\ref{fig:mmgbsa}). The strongest single contributor outside the pocket is CD4 Arg59 ($-5.98$~kcal/mol), a known interface anchor residue.

\begin{figure}[!ht]
\centering
\includegraphics[width=\textwidth]{F13_mmgbsa_decomposition.png}
\caption{\textbf{MM-GBSA per-residue energy decomposition of the relay pocket.} (A)~Per-residue binding free energy contributions (kcal/mol) of the eight relay pocket residues to the gp120--CD4 interaction in the reduced system. Gly467--Gly469 dominate, contributing $-3.40$~kcal/mol collectively. (B)~Overall energy component breakdown: van der Waals, electrostatic, polar solvation (EGB), and nonpolar solvation (ESURF) contributions to the total binding free energy ($-31.4$~kcal/mol).}
\label{fig:mmgbsa}
\end{figure}

\textbf{Drug binding free energy.} MM-GBSA on the 10~ns drug-bound ternary complex yields $\Delta G_\mathrm{bind} = -37.7 \pm 3.5$~kcal/mol for BMS-626529 at the relay pocket (40 frames), dominated by van der Waals packing ($-68.1$~kcal/mol). Per-residue decomposition within 6~\AA\ identifies 9 residues contributing $>1$~kcal/mol each: Leu115 ($-3.0$), Val208 ($-2.8$), Phe210 ($-2.6$), Val254 ($-1.4$), Lys116 ($-1.3$), Pro206 ($-1.3$), Pro212 ($-1.2$), Phe374 ($-1.0$), and Phe380 ($-1.0$~kcal/mol; PDB 6MEO chain~G numbering). All are \textbf{residues lining the relay pocket cavity} spanning gp120 domains C1, C2, and C3, confirming the ligand occupies the predicted C5 binding site. Interaction entropy analysis ($-T\Delta S_\mathrm{IE} = +17.67$~kcal/mol, corrected $\Delta G = -27.56$~kcal/mol) is reported in S6~Text; the high interaction energy variance ($\sigma = 14.40$~kcal/mol) exceeds the Duan~et~al.\ reliability threshold \citep{Duan2016}, so the enthalpy-dominated estimate ($-37.7$~kcal/mol) is the more robust metric.

\subsection{Cross-Structure Validation (SOSIP Trimers)}

Four additional HIV-1 Env structures representing different clades and conformational states were analyzed to assess relay pocket conservation beyond PDB 6MEO (Table~\ref{tab:crossval}).

\begin{table}[htbp]
\centering
\caption{Cross-structure validation of relay pocket across 5 HIV-1 Env structures.}
\label{tab:crossval}
\footnotesize
\begin{tabularx}{\textwidth}{@{}l >{\raggedright\arraybackslash}X ccccc@{}}
\toprule
PDB & Clade / Form & Backbone RMSD & GNM intra & fpocket & MD pocket RMSD & MD pocket RMSD \\
& & to 6MEO (\AA) & corr. & residues & (\AA, 20~ns) & (\AA, 500~ns) \\
\midrule
6MEO & BG505, CD4+CCR5 & -- & 0.69 & 7/8 & $0.71 \pm 0.15$\textsuperscript{a} & $1.97 \pm 0.40$\textsuperscript{b} \\
6MU6 & BG505, BMS-814508 & 0.72 & 0.69 & 4/8 & $1.34 \pm 0.27$\textsuperscript{b} & $\mathbf{1.72 \pm 0.44}$\textsuperscript{b} \\
5U7O & BG505, BMS-626529 & 0.78 & 0.73 & 8/8 & $1.07 \pm 0.23$\textsuperscript{b} & $1.79 \pm 0.56$\textsuperscript{b} \\
4TVP & BG505, prefusion & 0.54 & 0.73 & 5/8 & $1.29 \pm 0.34$\textsuperscript{b} & $\mathbf{1.98 \pm 0.44}$\textsuperscript{b} \\
5FYJ & Clade G, VRC01 (G471A) & 0.93 & 0.80 & 8/8 & $1.34 \pm 0.25$\textsuperscript{b} & $\mathbf{1.29 \pm 0.30}$\textsuperscript{b} \\
\bottomrule
\end{tabularx}

\vspace{0.5em}
{\footnotesize \textsuperscript{a}Full ternary complex (gp120+CD4+CCR5, 294,000 atoms). \textsuperscript{b}Reduced system (chain~G only, 85--101 residues, ${\sim}$33,000--38,000 atoms; no binding partners). All four SOSIP 500~ns pocket RMSDs are sub-2~\AA, with three of four \textit{lower} than the 6MEO 500~ns reference (1.97~\AA).}
\end{table}

All five structures show backbone RMSD $<1.0$~\AA\ at the relay segment, GNM intra-pocket correlation $\geq 0.69$, and pocket detection by at least one independent algorithm. Across all five structures, \textbf{28/28 GNM intra-pocket correlations are positive} (140/140 total), confirming that the relay segment moves as a cohesive mechanical unit regardless of conformational state. For the CD4-bound structure (6MEO), GNM analysis of the full ternary complex yields a relay--CD4 inter-domain coupling of 0.88 (perturbation response effectiveness: 98.3rd percentile), indicating direct mechanical communication between the relay pocket and the CD4-binding site. The SOSIP structures lack co-crystallized CD4, so inter-domain coupling was assessed only for 6MEO. GNM alanine scanning identified Asp470 as the hub residue in all five structures (11.8--18.7\% disruption upon virtual ablation), consistent with its central position in the relay backbone.

\textbf{SOSIP molecular dynamics cross-validation.} To extend the static cross-validation to dynamic assessment, 20~ns explicit-solvent MD simulations were performed for four SOSIP structures using the same protocol as the 6MEO reduced system (\S2.5): a 15~\AA\ shell around the relay pocket (complete residues, chain~G, 85--101 residues, ${\sim}$700--834 protein atoms), AMBER99SB-ILDN force field, TIP3P water, 0.15~M NaCl, ${\sim}$33,000--38,000 total atoms. The relay pocket maintained low internal RMSD in all four independent structures:

\begin{itemize}[nosep,leftmargin=1.5em]
    \item \textbf{6MU6} (2.55~\AA\ X-ray): pocket RMSD $1.34 \pm 0.27$~\AA\ (max 2.46~\AA), Rg $7.34 \pm 0.20$~\AA\ ($-1.9$\% change), $T = 300.0$~K
    \item \textbf{4TVP} (3.1~\AA\ X-ray): pocket RMSD $1.29 \pm 0.34$~\AA\ (max 2.42~\AA), Rg $7.39 \pm 0.23$~\AA\ ($+2.9$\% change), $T = 299.9$~K
    \item \textbf{5U7O} (3.0~\AA\ X-ray): pocket RMSD $1.07 \pm 0.23$~\AA\ (max 1.72~\AA), Rg $6.87 \pm 0.25$~\AA\ ($+4.9$\% change), $T = 300.0$~K
    \item \textbf{5FYJ} (3.1~\AA\ X-ray, G471A variant): pocket RMSD $1.34 \pm 0.25$~\AA\ (max 2.12~\AA), Rg $7.61 \pm 0.21$~\AA\ ($+6.2$\% change), $T = 300.0$~K
\end{itemize}

All four SOSIP trajectories show pocket RMSD of $1.07$--$1.34$~\AA, intermediate between the full ternary complex ($0.71$~\AA, with CD4+CCR5 stabilization) and the 500~ns 6MEO reduced system ($1.97$~\AA). The 5U7O structure shows the best pocket RMSD ($1.07$~\AA), and the 5FYJ G471A variant (RPAGGDMR) produces an identical pocket RMSD ($1.34$~\AA) despite the sequence change. These are single-chain reduced systems without binding partners, establishing that the relay pocket possesses intrinsic structural stability across independent structures, crystal forms, resolutions ($2.55$--$3.1$~\AA), and sequence variants.

\textbf{500~ns SOSIP extensions.} To assess pocket stability on physiologically relevant timescales, all four SOSIP structures were extended to 500~ns (50,001 frames each, 2,000~ns aggregate; Fig~\ref{fig:rmsd_timeline}). All four maintain sub-2~\AA\ pocket RMSD:

\begin{itemize}[nosep,leftmargin=1.5em]
    \item \textbf{5FYJ} (500~ns, G471A variant): pocket RMSD $1.29 \pm 0.30$~\AA, backbone RMSD $17.24 \pm 8.82$~\AA, RMSF 0.69~\AA\ (\textbf{0th percentile}), SASA $1{,}294 \pm 35$~\AA$^2$
    \item \textbf{6MU6} (500~ns): pocket RMSD $1.72 \pm 0.44$~\AA, backbone RMSD $6.90 \pm 1.01$~\AA, RMSF 1.19~\AA\ (\textbf{0th percentile}), SASA $1{,}279 \pm 33$~\AA$^2$
    \item \textbf{5U7O} (500~ns): pocket RMSD $1.79 \pm 0.56$~\AA, backbone RMSD $8.36 \pm 2.76$~\AA, RMSF 3.93~\AA\ (52nd percentile), SASA $1{,}212 \pm 55$~\AA$^2$
    \item \textbf{4TVP} (500~ns, unliganded): pocket RMSD $1.98 \pm 0.44$~\AA, backbone RMSD $7.20 \pm 1.75$~\AA, RMSF 1.32~\AA\ (\textbf{0th percentile}), SASA $1{,}254 \pm 46$~\AA$^2$
\end{itemize}

Three of four SOSIP structures show pocket RMSD \textit{lower} than the 6MEO 500~ns reduced system ($1.97$~\AA); 4TVP matches it ($1.98$~\AA). Three of four structures (5FYJ, 6MU6, 4TVP) operate at the 0th RMSF percentile (S3~Fig): no other residue in these reduced systems exhibits lower fluctuation over 500~ns. The relay pocket is a \textbf{conformational anchor} in the pre-triggered trimer across four independent X-ray structures ($2.55$--$3.1$~\AA\ resolution; S2~Fig), two HIV-1 clades (BG505 and X1193.c1), and both wildtype (RPGGGDMR) and variant (RPAGGDMR) sequences.

\begin{figure}[!ht]
\centering
\includegraphics[width=\textwidth]{F1_rmsd_timeline.png}
\caption{\textbf{Multi-structure relay pocket RMSD over 500~ns.} Pocket C$\alpha$ RMSD time series for five independent MD simulations: 6MEO reduced system, 5U7O SOSIP, 6MU6 SOSIP, 4TVP SOSIP, and 5FYJ G471A variant. All maintain sub-2~\AA\ pocket RMSD over 500~ns despite global backbone excursions up to 17~\AA. Dashed horizontal line at 2.0~\AA\ for reference. The relay pocket is kinematically decoupled from global tertiary fluctuations across all structures.}
\label{fig:rmsd_timeline}
\end{figure}

We define a structural element as \textit{kinematically decoupled} when its local RMSD remains below 30\% of the global backbone RMSD (pocket/backbone ratio $<0.30$), indicating that local stability is maintained independently of surrounding tertiary fluctuations. All four SOSIP systems satisfy this criterion: 5FYJ (0.075), 5U7O (0.214), 6MU6 (0.249), 4TVP (0.275). The strongest kinematic decoupling appears in 5FYJ: the pocket maintains $1.29$~\AA\ RMSD while the global backbone drifts to $17.24$~\AA\ (ratio 0.075). This 13-fold differential between local pocket stability and global backbone fluctuation is the most extreme in the dataset, confirming that the relay pocket's rigidity is an intrinsic property of the local sequence-structure relationship, independent of the surrounding tertiary context.

The increased pocket dynamics in the 6MEO ternary complex (RMSD $1.97$~\AA\ vs.\ $1.29$--$1.79$~\AA\ in three SOSIP structures) reflect the active propagation of the allosteric signal when CD4 and CCR5 are engaged. This observation is consistent with the FEP-measured $\Delta\Delta G = +3.95$~kcal/mol (\S3.15): the drug selectively stabilizes the hyper-rigid pre-triggered conformation, and the thermodynamic penalty quantifies the energy the conformational relay would need to overcome to reshape the pocket while the drug is bound.

Time-segmented analysis reveals distinct dynamic signatures. The 5U7O pocket exhibits conformational breathing (RMSD oscillates: $1.16 \to 1.92 \to 2.37 \to 1.52 \to 1.98$~\AA\ across five 100~ns segments), returning to low-RMSD states even as the global backbone drifts to 11.6~\AA. The 6MU6 pocket shows steady progression ($1.48 \to 1.64 \to 1.69 \to 1.71 \to 2.07$~\AA) with lower variance. The 5FYJ pocket (G471A) converges \textit{downward} ($1.37 \to 1.50 \to 1.29 \to 1.14 \to 1.17$~\AA), locking into its most rigid state by 300~ns despite the backbone undergoing the largest global excursion in the dataset ($17.24$~\AA). The 4TVP pocket (unliganded) shows a slow upward drift ($1.54 \to 1.69 \to 1.95 \to 2.24 \to 2.47$~\AA) consistent with the absence of stabilizing binding partners, yet remains sub-2.5~\AA\ throughout. In all four cases, the relay pocket is kinematically decoupled from global tertiary fluctuations.

\textbf{SOSIP mechanistic cross-validation (500~ns).} To extend cross-validation beyond structural stability to mechanistic characterization, dynamic cross-correlation matrices (DCCM), principal component analysis (PCA), and ensemble docking were performed on all four completed 500~ns SOSIP trajectories (5U7O, 6MU6, 5FYJ, and 4TVP) using protocols matching the 6MEO analyses (\S2.8, \S2.9).

\textit{Intra-pocket DCCM.} Covariance matrices were computed via \texttt{gmx covar} (All\_CA fit, dt=50~ps) and normalized to yield C$\alpha$--C$\alpha$ cross-correlations for the relay pocket (8 residues, 28 unique pairs; Fig~\ref{fig:dccm_comparison}).

\begin{table}[h]
\centering
\caption{Intra-pocket DCCM comparison across conformational states (500~ns).}
\label{tab:sosip_dccm}
\small
\begin{tabular}{@{}ll c c c c@{}}
\hline
System & Conformation & All 28 pos.? & Mean $C_{ij}$ & Seq.\ neigh.\ mean & Range \\
\hline
6MEO (500~ns) & CD4+CCR5 bound & Yes & $+0.79$ & $+0.97$ & $+0.37$ to $+0.99$ \\
6MU6 SOSIP (500~ns) & Pre-fusion (BG505) & No (24/28) & $+0.45$ & $+0.79$ & $-0.13$ to $+0.90$ \\
5U7O SOSIP (500~ns) & Pre-fusion (BG505) & No (19/28) & $+0.22$ & $+0.81$ & $-0.59$ to $+0.91$ \\
4TVP SOSIP (500~ns) & Pre-fusion (BG505) & No (20/28) & $+0.31$ & $+0.79$ & $-0.47$ to $+0.92$ \\
5FYJ SOSIP (500~ns) & G471A (Clade~G) & \textbf{Yes} & $+0.72$ & $+0.97$ & $+0.21$ to $+0.99$ \\
\hline
\end{tabular}
\end{table}

The SOSIP structures reveal a state-dependent DCCM signature absent from the 6MEO ternary complex. In all three wildtype SOSIP structures (6MU6, 5U7O, 4TVP), sequential neighbors remain strongly correlated ($+0.79$ to $+0.81$), but non-sequential pairs between the N-terminal (Arg465--Gly469) and C-terminal (Asp470--Arg472) halves of the pocket show anti-correlations: the two halves breathe in opposition. In the 6MEO ternary complex (CD4+CCR5 engaged), all 28 pairs are positive (mean $+0.79$), indicating that receptor engagement locks both halves into co-motion as the conformational relay signal propagates through the pocket.

The G471A variant (5FYJ) abolishes this gating motion. All 28 intra-pocket pairs are positive (mean $+0.72$), matching the 6MEO relay-active sign pattern despite the absence of CD4 or CCR5. The correlation signs match, but the underlying physics differ: 5FYJ has the lowest total PCA variance ($2.72$~nm$^2$) and lowest Schlitter entropy ($1{,}246$~J/mol$\cdot$K) of any structure in the dataset (Table~\ref{tab:sosip_pca}). The residues move together because they barely move at all. The anti-correlated breathing is gone because the alanine side chain at position 471 eliminates the backbone flexibility that the conserved glycine provides, converting a functional gate into steric gridlock. The wildtype Gly471 is required for the pocket to function as a gate (openable by the relay signal); the G471A variant seizes the gate shut.

The transition from anti-correlated breathing (wildtype SOSIP, pre-triggered) to fully correlated co-motion (6MEO, triggered) is reproduced across three independent wildtype structures from independent 500~ns simulations. This provides a direct mechanistic readout of the relay signal, consistent with the FEP-measured $\Delta\Delta G$ of $+3.95$~kcal/mol (\S3.15): a relay-blocking drug stabilizes the anti-correlated pre-triggered state, preventing the transition to correlated co-motion required for fusion activation. The G471A variant demonstrates that eliminating the gate has the same DCCM sign signature as activating the relay, but through rigidity rather than signal propagation.

\begin{figure}[!ht]
\centering
\includegraphics[width=\textwidth]{F5_dccm_comparison.png}
\caption{\textbf{State-dependent DCCM comparison across conformational states.} Intra-pocket dynamic cross-correlation matrices for (A)~6MEO ternary (CD4+CCR5 bound, relay active), (B)~5U7O SOSIP (wildtype pre-fusion), (C)~6MU6 SOSIP (wildtype pre-fusion), (D)~4TVP SOSIP (wildtype unliganded), and (E)~5FYJ SOSIP (G471A variant). Wildtype SOSIP structures show anti-correlated breathing between N- and C-terminal pocket halves (blue cells), while 6MEO and G471A show fully positive correlations (red cells) --- through active relay propagation and steric gridlock, respectively.}
\label{fig:dccm_comparison}
\end{figure}

\textit{PCA / essential dynamics.} Principal component analysis of the relay pocket atoms reveals structure-dependent conformational dynamics.

\begin{table}[h]
\centering
\caption{PCA of relay pocket dynamics across conformational states.}
\label{tab:sosip_pca}
\small
\begin{tabular}{@{}l c c c c@{}}
\hline
System & Total Var.\ (nm$^2$) & Top 3 PCs (\%) & Top 5 PCs (\%) & Schlitter $S$ (J/mol$\cdot$K) \\
\hline
5U7O SOSIP (500~ns) & 6.98 & 67.3 & 79.1 & 1472.6 \\
4TVP SOSIP (500~ns) & 5.44 & 61.7 & 71.7 & 1437.7 \\
6MU6 SOSIP (500~ns) & 4.45 & 59.1 & 71.2 & 1346.4 \\
5FYJ SOSIP (500~ns, G471A) & 2.72 & 54.2 & 63.9 & \textbf{1246.5} \\
6MEO reduced (100~ns) & 2.48 & 51.0 & 65.5 & -- \\
\hline
\end{tabular}
\end{table}

Total variance spans a 2.6-fold range across the four SOSIP structures ($2.72$--$6.98$~nm$^2$; S4~Fig), with conformational entropy following the same ordering. The G471A variant (5FYJ) has the lowest variance ($2.72$~nm$^2$) and lowest Schlitter entropy ($1{,}246$~J/mol$\cdot$K), consistent with the steric gridlock observed in DCCM: the alanine substitution restricts the backbone conformational space that glycine permits. Among wildtype structures, 6MU6 is the most rigid ($4.45$~nm$^2$, $1{,}346$~J/mol$\cdot$K), consistent with its 0th-percentile RMSF. In all five systems, dynamics are dominated by 2--3 collective modes, with the top 3 PCs capturing 51--67\% of variance, indicating a well-defined conformational landscape rather than stochastic fluctuation.

\textit{Ensemble docking.} GROMOS clustering (2.5~\AA\ cutoff on relay pocket C$\alpha$, dt=500~ps) of the 500~ns 5U7O and 6MU6 SOSIP trajectories yielded 3 cluster representatives each (reflecting the pocket's rigidity). BMS-626529 was docked into each representative using AutoDock Vina 1.2.7.

\begin{table}[h]
\centering
\caption{Ensemble docking of BMS-626529 across conformational states.}
\label{tab:sosip_ensemble}
\begin{tabular}{lcccc}
\hline
System & $N$ clusters & Mean affinity (kcal/mol) & Best & Worst \\
\hline
6MEO reduced (100~ns) & 20 & $-6.81 \pm 0.75$ & $-8.34$ & $-5.51$ \\
5U7O SOSIP (500~ns) & 3 & $-4.29 \pm 0.24$ & $-4.62$ & $-4.12$ \\
6MU6 SOSIP (500~ns) & 3 & $-3.90 \pm 0.33$ & $-4.37$ & $-3.64$ \\
\hline
\end{tabular}
\end{table}

The SOSIP ensemble docking scores ($-3.9$ to $-4.3$~kcal/mol; S5~Fig) are lower than the 6MEO ternary scores ($-6.8$~kcal/mol). This is expected because the SOSIP reduced systems lack the CD4 wall residues (Gln25, Asn39, Gln40, Phe43, Thr45) that form part of the binding pocket in the ternary complex. The ${\sim}2.5$~kcal/mol difference quantifies CD4's geometric contribution to the binding site: CD4 provides a physical wall that increases Van der Waals contacts and constrains ligand orientation, improving the enthalpy-dominated docking score. The SOSIP docking scores nonetheless remain in the druggable range for unoptimized probes ($>-3.5$~kcal/mol), demonstrating that the pocket possesses intrinsic druggability independent of CD4. This resolves an apparent paradox with the FEP results (\S3.15): while docking (enthalpy-only) favors the CD4-bound geometry, FEP (full thermodynamics including entropy) shows the drug binds preferentially to the pre-CD4 state ($\Delta G_\mathrm{bind} = -4.0$~kcal/mol pre-CD4 vs $-0.1$~kcal/mol post-CD4). The discrepancy reflects the large entropic penalty of forcing the drug into the rigidified, relay-activated ternary conformation, a penalty invisible to rigid-body docking but captured by alchemical free energy methods. The few SOSIP clusters (3 each, vs 20 for 6MEO) further confirm the pocket's exceptional rigidity in the pre-triggered state.

\subsection{Allosteric Network Analysis}

Dynamic cross-correlation analysis of the 500~ns trajectory (10,001 frames) reveals that the relay pocket moves as a highly cohesive mechanical unit. All 28 intra-pocket residue pairs show \textbf{positive cross-correlations} ($C_{ij} = +0.37$ to $+0.99$; mean $+0.79$), with sequential neighbors exhibiting near-unity coupling (mean $+0.97$) and a smooth distance-dependent decay. These correlations are substantially stronger than those observed in the initial 100~ns analysis ($C_{ij} = +0.08$ to $+0.82$), consistent with the improved statistical convergence afforded by five-fold greater sampling.

The relay pocket is thus identified as a \textbf{cohesive mechanical unit} in the gp120--CD4 allosteric network: high intra-pocket correlation across the full half-microsecond trajectory, combined with sub-2~\AA\ internal RMSD (\S3.5), indicates that the RPGGGDMR motif functions as a structurally rigid relay element.

Comparison with the SOSIP DCCM (\S3.11, Table~\ref{tab:sosip_dccm}) reveals a \textbf{state-dependent correlation signature}: in the three wildtype pre-fusion SOSIP trimers, intra-pocket correlations are weaker (mean $+0.22$ to $+0.45$) with anti-correlations between the N- and C-terminal halves of the motif, while in the 6MEO ternary complex (relay activated), all 28 pairs show strong positive correlation (mean $+0.79$). This transition from anti-correlated breathing to fully correlated co-motion, reproduced across three independent wildtype SOSIP structures from independent 500~ns simulations, is the dynamic fingerprint of the conformational relay signal. The G471A variant (5FYJ) produces the same all-positive DCCM sign pattern (28/28, mean $+0.72$) but through steric gridlock rather than active signal propagation, confirming that Gly471 is required for the gate to function.

\subsection{Allosteric Silencing by Drug Binding}

To directly test the relay blocker mechanism, DCCM was computed on the full ternary complex (887 C$\alpha$ atoms) from both the 20~ns apo and 10~ns drug-bound (BMS-626529 analog) trajectories, enabling comparison of cross-domain correlation patterns with and without a relay pocket--bound ligand.

In the apo ternary complex, the relay pocket residues (372--379 in GROMACS numbering; RPGGGDMR) exhibit near-perfect internal correlation ($\bar{C}_{ij} = 0.987$, 28/28 pairs positive, range 0.96--1.00) and maintain significant cross-domain coupling: relay$\to$CCR5 $|C| = 0.47$, relay$\to$CD4 $|C| = 0.49$. The relay pocket is dynamically coupled to both binding partners.

In the drug-bound system, a selective asymmetric silencing pattern emerges (Fig~\ref{fig:allosteric_silencing}):

\begin{itemize}[nosep,leftmargin=1.5em]
    \item \textbf{Relay$\to$CD4 coupling}: $|C| = 0.49 \to 0.24$ (\textbf{$-51\%$}): the incoming conformational signal from CD4 engagement is halved
    \item \textbf{Relay$\to$gp120 (intra-chain)}: $|C| = 0.84 \to 0.58$ ($-31\%$): internal signal propagation through gp120 is disrupted
    \item \textbf{Relay$\to$CCR5 coupling}: $|C| = 0.47 \to 0.40$ ($-15\%$): downstream coupling to the coreceptor is relatively preserved
    \item \textbf{Intra-pocket correlation}: $0.987 \to 0.894$ ($-9.4\%$): the pocket remains structurally intact but transmits less efficiently
    \item \textbf{Global mean $|C|$}: $0.519 \to 0.399$ ($-23\%$): system-wide decorrelation
\end{itemize}

The drug thus acts as a \textbf{conformational signal wedge}: the pocket retains its structural identity ($\bar{C}_{ij} = 0.89$) but can no longer receive and transmit the allosteric signal from CD4 engagement. The 3.4-fold selectivity in silencing (51\% CD4 vs.\ 15\% CCR5) is consistent with a relay blocker that intercepts the incoming conformational trigger from CD4 engagement, preventing the relay from transmitting the activation signal to the downstream fusion machinery. This pattern distinguishes relay blockers from attachment inhibitors (fostemsavir, which blocks CD4 binding) and coreceptor antagonists (maraviroc, which blocks CCR5 engagement).

\begin{figure}[!ht]
\centering
\includegraphics[width=\textwidth]{F11_allosteric_silencing.png}
\caption{\textbf{Allosteric silencing by drug binding.} (A)~Intra-pocket DCCM mean and range in apo versus drug-bound ternary complex, showing retained structural identity with reduced transmission. (B)~Cross-domain correlation changes upon BMS-626529 binding: relay$\to$CD4 coupling drops 51\%, relay$\to$gp120 drops 31\%, while relay$\to$CCR5 is relatively preserved ($-15$\%). The drug intercepts the incoming CD4 conformational signal at the relay junction, preventing downstream transmission to the fusion machinery.}
\label{fig:allosteric_silencing}
\end{figure}

\subsection{Non-Competitive Mechanism Validation}

To test whether the relay pocket--bound drug interferes with CD4 binding, MM-GBSA was computed for the gp120--CD4 interaction in the drug-bound ternary complex (BMS-626529 analog bound at the relay pocket) using 200 frames from the ligand-bound trajectory. The gp120--CD4 binding free energy in the drug-bound state is $\Delta G_\mathrm{CD4}^\mathrm{drug} = -28.39 \pm 8.94$~kcal/mol, compared to $-38.02 \pm 28.04$~kcal/mol in the apo ternary complex, yielding $\Delta\Delta G = +9.63$~kcal/mol. The CD4 interaction remains strongly favorable ($-28.4$~kcal/mol) despite the presence of BMS-626529 in the relay pocket, confirming that the drug does not competitively displace CD4. The $+9.63$~kcal/mol difference is modest relative to the magnitude of binding and falls within the standard deviations of both estimates.

Together with the allosteric silencing (\S3.13), this establishes a consistent picture: the drug permits CD4 engagement while disrupting the downstream conformational relay.

\subsection{Allosteric Coupling Free Energy (FEP/TI)}

To quantify the thermodynamic coupling between CD4 engagement and drug binding at the relay pocket, alchemical free energy perturbation was performed on three systems: a ternary complex (gp120{+}CD4{+}BMS-626529), a binary system (gp120{+}BMS-626529, no CD4), and a solvent reference (BMS-626529 in water/ions; Fig~\ref{fig:fep_cycle}). Three independent replicas with different velocity seeds were computed for each system (9 legs total, 1,440~ns aggregate; Table~\ref{tab:fep}).

\begin{table}[htbp]
\centering
\caption{FEP drug decoupling free energies (mean $\pm$ SEM across 3 replicas; 32 $\lambda$ windows $\times$ 5~ns each) and binding free energies via the standard thermodynamic cycle ($\Delta G_\mathrm{bind} = \Delta G_\mathrm{solvent} - \Delta G_\mathrm{complex}$).}
\label{tab:fep}
\small
\begin{tabular}{@{}lcccc@{}}
\toprule
& $\Delta G_\mathrm{ternary}$ & $\Delta G_\mathrm{binary}$ & $\Delta G_\mathrm{solvent}$ & $\Delta G_\mathrm{bind}$ \\
& (kcal/mol) & (kcal/mol) & (kcal/mol) & (kcal/mol) \\
\midrule
\multicolumn{5}{@{}l}{\textit{Decoupling free energies (higher = stronger binding):}} \\
TI$^\dagger$ & $6.63 \pm 0.01$ & $10.57 \pm 0.18$ & $6.54 \pm 0.01$ & -- \\
MBAR & $6.55 \pm 0.01$ & $10.37 \pm 0.19$ & $6.51 \pm 0.02$ & -- \\
\midrule
\multicolumn{5}{@{}l}{\textit{Binding free energies (negative = favorable):}} \\
& gp120{+}CD4 & gp120 only & & $\Delta\Delta G$ \\
TI$^\dagger$ & $-0.09 \pm 0.02$ & $-4.03 \pm 0.18$ & & $+3.95 \pm 0.18$ \\
MBAR & $-0.04 \pm 0.03$ & $-3.86 \pm 0.19$ & & $+3.82 \pm 0.19$ \\
\multicolumn{5}{@{}l}{\footnotesize $^\dagger$Per-window decorrelated before integration.} \\
\bottomrule
\end{tabular}
\end{table}

TI and MBAR estimators agree within 0.2~kcal/mol across all three legs, and inter-replicate standard deviations are small (ternary: 0.03, binary: 0.30--0.32, solvent: 0.02--0.04~kcal/mol), confirming convergence. The solvent leg serves as an internal consistency control: its inter-replicate SD ($<$0.04~kcal/mol) is an order of magnitude below the CD4 effect, confirming that the coupling signal is real rather than a sampling artifact.

Applying the standard thermodynamic cycle, BMS-626529 binds gp120 alone (binary, pre-CD4 conformation) with $\Delta G_\mathrm{bind} = -4.03 \pm 0.18$~kcal/mol (TI) / $-3.86 \pm 0.19$~kcal/mol (MBAR), representing moderate affinity consistent with an unoptimized probe compound at a novel site ($K_d \approx 1.4$~mM at 310~K), providing a starting point for structure-based optimization. In the CD4-bound ternary complex, where the conformational relay has propagated through gp120, binding drops to $\Delta G_\mathrm{bind} = -0.09 \pm 0.02$~kcal/mol (TI) / $-0.04 \pm 0.03$~kcal/mol (MBAR). The binary system uses the 6MEO gp120 backbone conformation with CD4 computationally removed; this represents the CD4-competent backbone absent the CD4 protein mass, not a crystallographically determined pre-triggered state. The measured $\Delta\Delta G$ therefore quantifies the thermodynamic coupling between CD4 presence and drug binding, establishing that the two binding events are energetically anti-correlated regardless of whether this coupling operates through direct steric or indirect allosteric mechanisms. The allosteric coupling between the two sites is $\Delta\Delta G = +3.95 \pm 0.18$~kcal/mol (TI) / $+3.82 \pm 0.19$~kcal/mol (MBAR). By Le Chatelier's principle, this coupling is reciprocal: just as CD4 engagement imposes a $+3.95$~kcal/mol penalty on drug binding, drug binding imposes an equivalent thermodynamic penalty on the CD4-induced conformational relay; the drug acts as a conformational lock on the pre-triggered state.

\begin{figure}[!ht]
\centering
\includegraphics[width=0.85\textwidth]{F7_fep_cycle.png}
\caption{\textbf{FEP thermodynamic cycle.} Schematic of the alchemical free energy perturbation thermodynamic cycle. BMS-626529 is decoupled from three environments: ternary complex (gp120+CD4), binary (gp120 only), and solvent. Binding free energy in the pre-CD4 state ($-4.0$~kcal/mol) drops to near-zero ($-0.1$~kcal/mol) upon CD4 engagement, yielding $\Delta\Delta G = +3.95$~kcal/mol --- the first thermodynamic measurement of the conformational relay signal strength.}
\label{fig:fep_cycle}
\end{figure}

This result has a direct mechanistic interpretation. The drug exhibits state-dependent affinity characteristic of conformational selection (analogous to Imatinib's selectivity for the inactive DFG-out kinase conformation; \citealt{Schindler2000}): BMS-626529 binds the pre-CD4 relay pocket conformation with moderate affinity ($\Delta G_\mathrm{bind} \approx -4$~kcal/mol), selectively stabilizing the pre-triggered state. The $+3.95$~kcal/mol $\Delta\Delta G$ quantifies the thermodynamic strength of the conformational relay signal: the free energy that CD4-induced rearrangement would need to overcome to reshape the pocket while the drug is bound. A relay blocker drug exploits this coupling: by occupying the pre-CD4 pocket, it imposes a ${\sim}4$~kcal/mol thermodynamic barrier against the conformational relay, effectively locking gp120 in the pre-triggered conformation and preventing downstream fusion activation.

Together with the DCCM allosteric silencing (\S3.13) and MM-GBSA non-competitiveness analysis (\S3.14), this FEP result provides a fourth independent line of evidence for the relay blocker mechanism. The FEP and MM-GBSA measure reciprocal sides of the same coupling: MM-GBSA quantifies the drug's effect on CD4 binding ($\Delta\Delta G_\mathrm{CD4} = +9.63$~kcal/mol); FEP quantifies CD4's effect on drug binding ($\Delta\Delta G_\mathrm{drug} = +3.9$~kcal/mol). Both agree qualitatively: the relay pocket and CD4 binding site are thermodynamically coupled, with perturbation of one site propagating to the other. The quantitative difference reflects known systematic differences between endpoint and alchemical methods \citep{Genheden2015} and the presence of CCR5 in the MM-GBSA system but not in the FEP systems.

\subsection{Resistance Mutation Mapping}

Cross-referencing the relay pocket against the Stanford HIV Drug Resistance Database (28 known resistance-associated mutations) reveals that \textbf{zero known drug resistance mutations directly affect relay pocket residues 465--472}. The nearest resistance mutation is M475I (a secondary fostemsavir escape mutation), located 3 residues from the pocket boundary. Two clade-specific polymorphisms exist within the pocket (R472K, D470N), but both are conservative substitutions that preserve the backbone geometry and relay function.

A comprehensive PDB survey of 531 HIV-1 gp120-containing entries identified 360 structures with the relay pocket region resolved, of which \textbf{all 360 (100\%) contain the exact RPGGGDMR sequence}: zero variants across 27 years of structures (1998--2025), 115 X-ray and 245 cryo-EM, from multiple independent laboratories.

\subsection{Cross-Target Specificity}

To assess whether the relay pocket motif could produce off-target binding to host proteins, the complete RPGGGDMR octapeptide was searched against the entire reviewed human proteome (Swiss-Prot, 20,431 proteins). \textbf{Zero matches were found}. Relaxing the search to any 7-residue sub-motif (RPGGGDM, PGGGDMR) also yielded zero matches. The shortest sub-motif with any human proteome occurrence is the 5-mer RPGGG (30 proteins, 0.15\%) and PGGGD (31 proteins, 0.15\%), both scattered across unrelated proteins with no structural or functional relationship to gp120. The relay pocket's defining backbone geometry is completely absent from the human proteome, indicating negligible risk of off-target geometric mimicry for a drug designed against this motif.

\subsection{Computational Alanine Scanning}

Contact-weighted GNM alanine scanning of all 8 relay residues (S7~Fig) reveals that the relay signal propagates through backbone topology, not side-chain contacts: the glycine triplet (Gly467--469) shows \textbf{zero coupling change upon mutation to alanine} (glycine already has the minimal side chain), and the maximum coupling disruption from any single mutation is only $-2.2$\% (Arg465Ala). BMS-626529 re-docking against each mutant receptor shows preserved or improved pocket geometry (Asp470Ala: $-8.1$~kcal/mol; Arg472Ala: $-8.7$~kcal/mol vs.\ wild-type $-4.2$~kcal/mol), while vicriviroc binding is uniformly weakened ($+0.9$ to $+1.8$~kcal/mol). Full per-mutant results are in S6~Text.

\subsection{Ensemble Docking}

Docking of BMS-626529 across 20 representative MD conformations (GROMOS clustering, 2.5~\AA\ cutoff from the first 100~ns of the trajectory) yielded a mean binding affinity of $-6.81 \pm 0.75$~kcal/mol (range $-5.51$ to $-8.34$~kcal/mol). The ensemble mean ($-6.81$) is comparable to the single cryo-EM frame result ($-6.7$), confirming that pocket druggability is \textbf{maintained across the MD conformational ensemble}. Vicriviroc (a CCR5 antagonist, as a selectivity control) yielded $-7.50 \pm 0.82$~kcal/mol across the same 20 frames. The comparable docking scores are expected: AutoDock Vina's scoring function correlates strongly with heavy-atom contact count, and both drug-sized molecules (29 and 36 heavy atoms, respectively) can exploit the pocket's large bounding volume (1,588~\AA$^3$). Docking validates druggability (can a drug-sized molecule achieve favorable contacts?) rather than selectivity (does the pocket discriminate between scaffolds?), and the comparable scores for two structurally unrelated compounds independently confirm that the pocket geometry is compatible with drug-like binding. The best-scoring conformation (cluster~5) yielded $-8.34$~kcal/mol for BMS-626529 and $-9.06$~kcal/mol for vicriviroc, representing an optimal pocket geometry that could guide pharmacophore refinement.

\subsection{Hydration Site Analysis, PCA, and De Novo Design}

Hydration site analysis identified 22 water positions at $>$1.5$\times$ bulk density near the pocket, none of which were structural waters; eight displaceable sites (1.7--2.3$\times$ bulk) indicate that ligand binding would not require displacing tightly bound waters \citep{Halgren2009} (S6~Text). PCA of the first 100~ns (10,110 frames, 112 pocket atoms) shows three dominant collective modes capturing 51.0\% of variance (S4~Fig), consistent with structured dynamics rather than random fluctuation. De novo ligand design against the 55-feature pharmacophore yielded eight novel scaffolds (all Lipinski/Veber/PAINS-compliant; five with Tanimoto $<0.2$ to any known compound), demonstrating that the pocket geometry accommodates diverse chemical series (\texttt{denovo\_designs.json}, S6~Text).

\section{Discussion}

\subsection{A Gap in the Entry Inhibitor Landscape}

Four approved drugs target three of four distinct steps in the HIV-1 entry cascade (\S1), yet the conformational relay between receptor engagement and fusion activation has never been targeted with a small molecule. The gp120 Arg465--Arg472 pocket fills this gap.

Every existing entry inhibitor blocks a step. This pocket does not block a step. It terminates a signal. The conformational signal from CD4 to gp41 routes through this eight-residue junction. A drug in the pocket lets the signal arrive and terminates it. The virus binds its receptor, engages its coreceptor, and the conformational change required to activate fusion dissipates at the relay. Half-microsecond MD simulations of four independent SOSIP trimers (2,000~ns aggregate; \S3.11, Fig~\ref{fig:rmsd_timeline}) confirm that the relay pocket functions as a \textbf{conformational anchor} within the pre-triggered Env trimer: all four structures maintain sub-2~\AA\ pocket RMSD ($1.29$--$1.98$~\AA), three of four operate at the 0th RMSF percentile, and the pocket is kinematically decoupled from global backbone dynamics (ratio as low as 0.075 in 5FYJ, where the pocket holds $1.29$~\AA\ while the backbone drifts to $17.24$~\AA). This rigidity is an intrinsic property of the local sequence-structure relationship, reproduced across two HIV-1 clades (BG505 and X1193.c1), four independent X-ray structures ($2.55$--$3.1$~\AA), and both wildtype and variant (G471A) sequences. Pocket dynamics increase upon CD4 engagement: the 6MEO ternary complex pocket RMSD rises to $1.97$~\AA, reflecting active propagation of the allosteric signal as the conformational change occurs. A drug that locks the pocket in its hyper-rigid pre-triggered conformation blocks the conformational relay at its most mechanically constrained point.

This rigidity signature is mechanistically cross-validated by SOSIP DCCM analysis (\S3.11, Fig~\ref{fig:dccm_comparison}, Table~\ref{tab:sosip_dccm}): in all three wildtype pre-fusion SOSIP trimers (6MU6, 5U7O, 4TVP), the N- and C-terminal halves of the pocket exhibit anti-correlated breathing (mean $C_{ij} = +0.22$ to $+0.45$, with 4--9 negative pairs per structure), while in the 6MEO ternary complex (relay activated), all 28 intra-pocket pairs are positively correlated (mean $+0.79$). This state-dependent transition is reproduced across three independent wildtype structures from separate 500~ns simulations. The G471A variant (5FYJ) provides a natural control: it abolishes the anti-correlated breathing (28/28 positive, mean $+0.72$) but through steric gridlock, not signal propagation (lowest PCA variance, lowest Schlitter entropy in the dataset). Glycine-471 is the gate residue: the wildtype pocket breathes; the G471A pocket is seized shut. A drug binding the wildtype pre-triggered pocket stabilizes the anti-correlated resting state, preventing the transition to correlated co-motion required for downstream signal propagation.

The mechanism of inhibition is steric blockade of relay motion: a small molecule occupying this pocket permits normal CD4 binding (the pocket is approximately 15~\AA\ from the CD4 binding site) and normal CCR5 engagement (different protein surface), but disrupts the conformational relay at the gp120 C-terminal junction by constraining the rearrangement required for downstream fusion activation. The virus docks to the target cell but fails to trigger membrane fusion. Four independent computational lines of evidence support this non-competitive relay blocker mechanism: (1)~DCCM allosteric silencing shows the drug reduces relay$\to$CD4 coupling by 51\% while preserving relay$\to$CCR5 coupling ($-15$\% change; \S3.13, Fig~\ref{fig:allosteric_silencing}); (2)~MM-GBSA non-competitiveness analysis confirms that CD4 binding remains strongly favorable ($-28.4$~kcal/mol) in the drug-bound state (\S3.14); (3)~intra-pocket correlation analysis shows the relay pocket retains its structural identity ($\bar{C}_{ij} = 0.89$) but loses signal transmission capacity; and (4)~alchemical free energy perturbation (9 independent legs across 3 replicas) reveals strong negative allosteric coupling between the drug and CD4 ($\Delta G_\mathrm{bind}$: $-4.0$~kcal/mol pre-CD4 $\to$ $-0.1$~kcal/mol post-CD4; $\Delta\Delta G = +3.9$~kcal/mol; \S3.15, Fig~\ref{fig:fep_cycle}), providing the first thermodynamic measurement of the conformational relay signal strength and establishing that the drug selectively stabilizes the pre-triggered conformation. The FEP establishes that more than half the binding mechanism is not molecular contacts but the drug selecting a conformation: $+3.95$~kcal/mol is the energy that CD4-induced rearrangement transmits through the relay junction, providing the first direct measurement of the conformational relay signal strength. Experimental validation of this mechanism is outlined in \S4.8.

Independent biophysical data confirm the state-dependent dynamics predicted by our DCCM analysis. Hydrogen-deuterium exchange mass spectrometry (HDX-MS) of gp120 shows that the relay pocket region (peptic Fragment~15: T455--L483, spanning all of RPGGGDMR) exhibits moderate backbone protection in the unliganded state and becomes ``significantly stabilized'' upon CD4 binding \citep{Kong2010}, while full-length glycosylated gp120 shows ``moderate protection'' at residues 458--474 that increases to ``structural rigidification upon sCD4 binding'' \citep{Guttman2012}. This transition from dynamic (pre-CD4) to rigid (post-CD4) independently corroborates the relay gate mechanism: our DCCM shows anti-correlated breathing in pre-triggered SOSIP and all-positive co-motion after relay activation. Single-molecule FRET studies demonstrate three conformational states (pre-triggered, intermediate, and post-CD4; \citealt{Ma2018}), and gp41-directed bNAbs allosterically remodel gp120 dynamics through the same relay junction \citep{Xu2026}, providing direct experimental evidence for the bidirectional allosteric pathway our relay mechanism predicts. NMR confirms that the C5 domain immediately downstream of the relay pocket binds the gp41 ectodomain ($K_d \approx 1$~$\mu$M; \citealt{Guilhaudis2002}), establishing the structural basis for signal transmission from gp120 to gp41. This converges with network-theoretic analyses showing that allosteric communication pathways in gp120 connect CD4-binding-site rearrangements to distant functional regions \citep{Henderson2020}.

\subsection{Conservation and Resistance Barrier}

The relay segment (mean entropy 0.21 across 981 NCBI gp120 sequences) sits within the C5 subdomain (mean entropy 0.14), the most conserved constant region of gp120. Three residues (Arg465, Pro466, and Gly468) are completely invariant, and three more (Gly469, Met471, Arg472) exceed 94\% conservation. The RPGGGDMR motif is conserved at $>$96\% identity across all major HIV-1 Group~M clades (A through G and circulating recombinant forms); 766 of 981 sequences carry the exact motif. Even the most variable position (Gly467, entropy 0.72) shows predominantly G$\to$A substitutions that preserve the small side-chain character required for relay backbone flexibility.

The structural basis for this constraint: Pro466 introduces a backbone kink required for relay geometry and is completely invariant. Of the triple glycine (467--469), Gly468 is invariant and Gly469 is 98.5\% conserved; Gly467 tolerates alanine substitution (9\% of sequences), consistent with a requirement for small side chains rather than absolute glycine identity at this position. Asp470 shows D$\to$N substitutions (8\%) that preserve the polar character of the charge anchor. Computational alanine scanning (\S3.18) confirms that the relay signal propagates through backbone topology, not side chains, with the glycine triplet showing zero coupling change upon mutation. The relay backbone is structurally constrained against mutation.

The G471A variant (PDB 5FYJ, Clade~G) provides direct dynamic evidence for this constraint. In 500~ns MD, the G471A pocket is the most rigid in the dataset (RMSD $1.29$~\AA, RMSF 0th percentile, Schlitter entropy $1{,}246$~J/mol$\cdot$K), and its DCCM shows all 28 intra-pocket pairs positive (mean $+0.72$), abolishing the anti-correlated breathing that all three wildtype SOSIP structures exhibit (\S3.11, Table~\ref{tab:sosip_dccm}). The alanine substitution eliminates the backbone flexibility that glycine provides at position 471, converting the relay gate from a functional switch (openable by CD4-induced signal) to a seized mechanism. Greater pocket rigidity is not an advantage; it is a loss of function. The virus conserves Gly471 because the relay must be a gate, not a wall.

Deep mutational scanning (DMS) of BG505.T332N SOSIP.664 \citep{Haddox2018} provides independent functional evidence for the relay gate. Five of eight relay pocket residues (Arg465, Pro466, Gly468, Gly469, Asp470) are rank~1 in single-round infectivity (the most fitness-constrained rank), confirming extreme functional conservation at these positions. Three positions show fitness landscape gaps: Met471 (rank~3/20) and Arg472 (rank~3/20) are naturally conserved at $>$97\% yet moderately DMS-tolerant, while the most dramatic outlier is position 471 (BG505 numbering; Gly467 in HXB2), where glycine ranks only 6th/20 and alanine ranks 5th/20 in DMS, yet natural conservation is 87\% glycine. This ``fitness landscape gap'' --- DMS-tolerant yet evolutionarily constrained --- is the signature of a functional role that extends beyond single-round entry. The relay gate model explains the Gly467 pattern: G471A preserves entry in a single-round assay (which does not require sustained relay signaling) but abolishes the backbone flexibility required for multi-round replication fitness where the conformational relay must operate reliably across many entry events. Evolutionary coevolution analysis (353 HIV-1 Env sequences, mutual information with average product correction) reveals that the two variable relay positions (Asp470 and Arg472) co-evolve as a structural pair (MI~$= 0.443$), consistent with a salt bridge that stabilizes the motif fold, and that Arg472 shows coevolution coupling to Lys486 in the gp41-contact region (MI~$= 0.393$), providing independent evolutionary evidence for the gp120-to-gp41 relay axis.

Five independent methods converge on a single gate-hinge residue within the relay motif: DMS (Gly467/471 rank~6/20, most DMS-tolerant position), AlphaFold (pLDDT 89.6, lowest in motif), Shannon entropy (0.72, highest in motif), cryo-EM local density (27th percentile, lowest in motif), and GNM normal modes (9th percentile mobility). This multi-method convergence on a single hinge residue, using fundamentally different data sources (fitness, structure prediction, sequence conservation, experimental density, and elastic network mechanics), constitutes strong independent evidence that the relay pocket contains a functional gate rather than a randomly flexible region.

This is further supported by resistance mutation mapping against the Stanford HIV Drug Resistance Database: \textbf{zero of 28 known drug resistance mutations affect relay pocket residues 465--472} (\S3.16). The nearest resistance mutation (M475I, a secondary fostemsavir escape) is 3 residues away. The PDB-wide structural survey (\S3.16) confirms: 360/360 solved gp120 structures carry the exact RPGGGDMR sequence. The virus has evolved resistance to every class of drug at other sites. It has never mutated this pocket. The relay is load-bearing: changing it to escape a drug would break the entry mechanism itself.

\subsection{Comparison to Nearest Prior Art}

The closest existing drug target to the relay pocket is fostemsavir's binding site on the gp120 $\beta$20-$\beta$21 element \citep{Pancera2017}. While both are on gp120, they are approximately 15~\AA\ apart, involve different residues, and operate through fundamentally different mechanisms: fostemsavir prevents CD4 attachment entirely, while a relay blocker permits attachment and coreceptor engagement but disrupts the downstream conformational relay required for fusion activation.

\subsection{The CCR5 ECL3 Site}

The identification of CCR5 ECL3 as a gp120 contact surface is independently significant. Most structural and pharmacological work on CCR5 for HIV has focused on the N-terminus, ECL2, and the transmembrane helical bundle (where maraviroc binds). ECL3 has not been extensively characterized as an HIV-relevant interface. If validated experimentally, this site could represent an additional allosteric target on CCR5 distinct from the maraviroc pocket.

\subsection{The Glycan Shield}

HIV-1 gp120 is heavily glycosylated, with N-linked glycans shielding much of the protein surface from antibody and small-molecule access. Structural analysis of glycosylation sequons (NxS/T, x$\neq$P) across six Env structures reveals that the relay pocket occupies a relatively glycan-sparse zone. In PDB 6MEO (BaL strain), the nearest sequon to the pocket centroid is N384 at 15.7~\AA, followed by N276 (15.8~\AA) and N262 (16.1~\AA); no resolved sugar atoms fall within 15~\AA\ of the pocket. In BG505 SOSIP structures (5U7O, 6MU6, 4TVP), the nearest sequon is N363 at ${\sim}$7.8~\AA\ (a site absent in the BaL strain, S363), yet even these structures contain at most one resolved sugar residue (a single NAG stub) within 15~\AA\ of the pocket center. PDB 5FYJ, the most densely glycosylated structure in our analysis (162 resolved sugar residues), contains zero sugars within 15~\AA\ of the pocket despite the bNAb PGT151 contacting this region \citep{Lee2016}. The glycan shield mapping of Stewart-Jones et~al.\ confirms that the C5 domain lies in a glycan-sparse corridor relative to the V1/V2 and outer-domain glycan clusters \citep{StewartJones2016}. The only FDA-approved small molecule targeting gp120, fostemsavir (BMS-626529), binds the $\beta$20--$\beta$21 element at a comparable distance from the glycan canopy \citep{Pancera2017}, demonstrating that small-molecule access to gp120 pockets is pharmacologically achievable despite the glycan shield. Small molecules (300--600~Da) require substantially less approach clearance than antibodies (${\sim}$150~kDa), and the pocket must remain conformationally accessible to execute its relay function, together suggesting favorable steric access for drug candidates.

\subsection{The Kinetic Window}

PDB 6MEO captures the ternary complex (gp120--CD4--CCR5), a transient state during viral entry. A key question is whether the relay pocket exists long enough for a drug to bind. Several considerations suggest the kinetic window is pharmacologically viable:

\begin{enumerate}[nosep,leftmargin=1.5em]
    \item \textbf{The drug's preferred state is pre-CD4.} FEP calculations (\S3.15) establish that the relay pocket binds BMS-626529 with moderate affinity ($-4.0$~kcal/mol) in the pre-CD4 conformation, with affinity dropping to near-zero upon CD4 engagement. This state-dependent binding profile indicates that the drug functions by conformational selection: it binds to and stabilizes the resting (pre-triggered) relay conformation before CD4 arrives. The kinetic window is therefore the entire prefusion state, not a brief post-CD4 intermediate.

    \item \textbf{Precedent for pre-binding.} Fostemsavir binds the prefusion trimer and remains bound through conformational changes. A relay-pocket drug similarly pre-loads, with its ${\sim}4$~kcal/mol thermodynamic penalty ($\Delta\Delta G$) resisting the CD4-induced conformational relay.

    \item \textbf{Structural rearrangement occurs at the relay.} The gp120 C-terminal region (including residues 465--472) undergoes conformational change upon CD4 binding \citep{Pancera2010,Herschhorn2016}. Pharmacophore analysis of the ternary complex reveals that the pocket wall includes residues from both gp120 (chain~G) and CD4 (chain~A: Gln25, Asn39, Gln40, Phe43, Thr45), indicating that CD4 residues approach the drug binding site upon engagement. If the drug-bound pre-CD4 conformation traps the pocket in a partially rearranged state, the drug may make contacts with both viral and host protein surfaces in the intermediate, creating an additional barrier to escape mutation since the virus cannot mutate the host half of the binding interface.

    \item \textbf{The rate-limiting step is CCR5 recruitment.} After CD4 binding, the virus must recruit CCR5 through V3 loop insertion, a slow, diffusion-limited process. A drug that locks the relay in the pre-triggered conformation would have this entire window to maintain its inhibitory effect.
\end{enumerate}

\subsection{Limitations}

All findings in this study are computational, derived from structures at 2.55--3.9~\AA\ resolution (cryo-EM and X-ray crystallography). Key limitations:

\begin{enumerate}[nosep,leftmargin=1.5em]
    \item \textbf{No experimental validation.} All evidence for the relay pocket is computational. The pocket's biological druggability, its accessibility in the native viral spike, and its capacity to block fusion when occupied have not been confirmed by biophysical or cellular assays. Specific experimental approaches to validate these predictions are detailed in \S4.8.

    \item \textbf{Glycan shield.} All simulations used bare (non-glycosylated) systems. Structural analysis of six Env structures shows zero resolved sugars within 15~\AA\ of the pocket center, and fostemsavir provides clinical precedent for small-molecule access to gp120 at comparable glycan proximity. However, the nearest glycosylation sequon is strain-dependent (N384 at 15.7~\AA\ in BaL versus N363 at ${\sim}$7.8~\AA\ in BG505 SOSIP), and dynamic glycan coverage cannot be assessed from static coordinates. Fully glycosylated Env trimer simulations at 500~ns--1~$\mu$s timescales are needed to map the dynamic approach corridor.

    \item \textbf{Structural resolution.} The primary structure (6MEO, 3.9~\AA\ cryo-EM) determines backbone positions but side-chain rotamers, particularly Arg465 and Arg472, are modeled rather than directly resolved. Cross-validation against a 2.55~\AA\ X-ray structure (6MU6) independently resolves the same pocket geometry including side-chain positions, and 500~ns MD confirms sub-angstrom stability (0th percentile RMSF, \S3.11). The pocket is reproduced in four additional structures spanning 2.55--3.1~\AA\ (\S3.3). The terminal guanidinium groups of Arg465 and Arg472 possess high intrinsic flexibility that static crystallography cannot fully constrain.

    \item \textbf{Conformational sampling.} Nine independent MD simulations (${\sim}4{,}000$~ns aggregate) across multiple starting structures confirm pocket persistence (RMSD 0.71--1.98~\AA). SOSIP mechanistic cross-validation (DCCM, PCA, ensemble docking on four 500~ns trajectories, 2,000~ns aggregate; \S3.11) extends validation to allosteric network properties and identifies the G471A gating mechanism. However, all trajectories sample microsecond timescales; millisecond-scale conformational transitions, including rare large-amplitude breathing, are not captured. Drug-bound SOSIP simulations and MM-GBSA on SOSIP systems remain as future work (\S4.1). Steered MD and umbrella sampling of the drug extraction coordinate are planned as future work to characterize the unbinding energy landscape and kinetic residence time.

    \item \textbf{Binding free energy decomposition.} MM-GBSA yields overall binding free energies of $-38.0$ and $-31.4$~kcal/mol for the ternary and reduced systems respectively, but per-residue energy decomposition was only achievable on the smaller reduced system (1,715 atoms). On the 14,149-atom ternary complex, all three decomposition modes (idecomp=1, 2, 3) failed to produce physically meaningful per-residue attributions, a known limitation of GB-based solvation energy decomposition on large multi-chain systems. Because the ternary context strengthens overall binding by ${\sim}6.6$~kcal/mol, the reduced-system pocket attribution ($-3.07$~kcal/mol) represents a truncated estimate of the true interaction energy.
\end{enumerate}

\subsection{Suggested Experimental Validation}

\textbf{For the gp120 relay pocket (Arg465--Arg472):}

\begin{enumerate}[nosep,leftmargin=1.5em]
    \item \textit{Alanine scanning mutagenesis.} Express Env trimers with individual alanine substitutions at positions 465--472. Measure CD4 and CCR5 binding (SPR or BLI), membrane fusion (cell--cell fusion assay), and pseudovirus entry (single-round infectivity assay). Expected: binding unaffected, fusion impaired.

    \item \textit{Hydrogen-deuterium exchange mass spectrometry (HDX-MS).} Measure deuterium incorporation in the gp120 C-terminal region during CD4-triggered conformational change. Published HDX-MS data show increased backbone protection at this region upon CD4 binding \citep{Kong2010,Guttman2012}; the key experiment is to test whether a relay-blocking compound restores pre-CD4 exchange levels (increased exchange relative to CD4-bound, decreased relative to unliganded), confirming that the drug locks the relay in its pre-triggered state.

    \item \textit{Virtual screening.} Screen ZINC, Enamine REAL, or other drug-like compound libraries against the relay pocket using the provided docking configurations. Top-scoring compounds can be tested in pseudovirus entry assays.
\end{enumerate}

\textbf{For CCR5 ECL3 (Tyr251, Leu255, Asn258, Asp276):}

\begin{enumerate}[nosep,leftmargin=1.5em]
    \item \textit{Alanine scanning mutagenesis.} Measure effect on gp120 binding affinity and HIV-1 entry.
    \item \textit{Cross-linking mass spectrometry (XL-MS).} Confirm direct contact between ECL3 residues and gp120 in the ternary complex.
\end{enumerate}

\section{Data Availability}

A comprehensive supplementary data package accompanies this study, containing 78 structured data files (JSON, XVG, DAT), 93 coordinate and docking pose files (PDB, PDBQT), 103 standalone analysis scripts (90 Python, 13 Bash), and full GROMACS simulation parameters (topology, MDP, and coordinate files). The package is organized into four top-level directories: \texttt{results/} (11 categories from pocket discovery through drug design, including FEP parameter files), \texttt{scripts/} (6 workflow categories from setup through drug design), \texttt{topology/} (force field parameters, MDP files, and run input), and \texttt{structures/} (PDB and PDBQT coordinate files for all receptors and ligands). All scripts use a relative path convention (\texttt{DATA\_ROOT}) that resolves to the package root directory, enabling direct execution without path modification. SHA-256 checksums are provided for all files. Raw MD trajectories for all five 500~ns simulations (${\sim}$25~GB) and FEP dH/d$\lambda$ data for all 288 lambda windows (${\sim}$2.5~GB) are available for download at \url{https://hiv1-relay-data.jretting.workers.dev} (self-contained packages with SHA-256 verification).

The data files span the full analytical scope of this work: 55-feature pharmacophore model, fpocket and P2Rank cavity detection, B-factor and druggability scoring, Shannon entropy from 981 NCBI gp120 sequences, AutoDock Vina docking for relay and control sites, 96-compound virtual screening results, ADMET predictions for 98 compounds, resistance mutation mapping (28 mutations), cross-clade structural survey (531 PDB entries, 360/360 exact RPGGGDMR), de novo lead designs, 20-frame ensemble docking, WaterMap hydration analysis, allosteric network and betweenness centrality, PCA essential dynamics, alanine scanning (8 mutants), proteome specificity (20,431 human proteins), SOSIP MD cross-validation (4 structures at 20~ns, all 4 extended to 500~ns with comprehensive analysis JSONs, 2,000~ns aggregate), SOSIP mechanistic cross-validation (intra-pocket DCCM and PCA with Schlitter entropy on all four 500~ns SOSIP trajectories including the G471A variant, ensemble docking on 5U7O and 6MU6), 500~ns DCCM (10,001 frames), MM-GBSA binding free energies (reduced and ternary systems, per-residue decomposition), IE-MMPBSA interaction entropy, non-competitiveness free energy analysis, full-ternary allosteric silencing DCCM (apo vs.\ drug-bound, 887 C$\alpha$ atoms), and FEP/TI alchemical free energies with TI, BAR, and MBAR estimators across replicate thermodynamic cycles.

The docking input files (ligand and receptor PDBQT files, 20 cluster frame conformations from GROMOS clustering) and all 40 ensemble docking output poses (20 BMS-626529 + 20 Vicriviroc) are included. FEP MDP parameter files are provided for independent binding free energy verification. The analysis scripts are standalone and self-documenting, requiring only publicly available software (GROMACS, AutoDock Vina, RDKit, MDAnalysis, ProDy, alchemlyb; see supplementary \texttt{README.md} for full dependency and execution documentation). Every quantitative claim in this paper can be independently verified by re-running the corresponding script against the provided data.

The complete supplementary data package (scripts, results, figures, structures, and topology files) is available at \url{https://github.com/jonretting/gp120-relay-supplementary}. A permanent archived version is deposited on Zenodo (DOI: \href{https://doi.org/10.5281/zenodo.23449544}{10.5281/zenodo.23449544}). The source structure PDB 6MEO is publicly available from the RCSB Protein Data Bank.

\section{Patent Disclosure}

A US provisional patent application covering the gp120 C5 conformational relay target and associated pharmacophore models has been filed (July 2026).

\section{AI Disclosure}

AI tools (Claude, Anthropic; Gemini, Google) were used for writing assistance, literature search, analysis script drafting, and scientific discussion including evaluation of mechanistic hypotheses and alternative interpretations. All molecular dynamics simulations and free energy perturbation calculations were performed by the author on self-owned hardware using GROMACS 2025.2 (see Data Availability). The evidentiary data underlying this paper---trajectories, binding free energies, cross-correlation matrices---are numerical solutions to classical equations of motion, not language model outputs; no AI system can generate molecular dynamics trajectory data. All 103 analysis scripts were validated via cold-start Docker testing (Python 3.11-slim, pip-only dependencies), producing identical quantitative results without any AI involvement. Every quantitative claim in this paper can be independently reproduced by executing the provided scripts against the provided simulation data. All figures were generated from raw simulation data by provided scripts. The author designed all experiments, selected all simulation methods and parameters, executed all calculations, and can independently explain and defend every methodological choice and result interpretation reported in this work.

\section*{Acknowledgments}

The author thanks Prof.\ David Tholen for correspondence that sharpened the author's approach to methodological rigor and critical evaluation of results.

\section*{Supporting Information}

\paragraph{S1 Fig. B-factor comparison across gp120 domains.} Mean experimental B-factors (PDB 6MEO) for the relay pocket compared to other gp120 regions (V3 loop, C2/bridging sheet, C4 domain) and binding partners (CD4, CCR5). The relay pocket (72.8~\AA$^2$, 38th percentile) is more ordered than V3, C2, and C4 domains.

\paragraph{S2 Fig. Cross-structure pocket overlay.} Structural superposition of the relay pocket backbone from five HIV-1 Env structures (6MEO, 5U7O, 6MU6, 4TVP, 5FYJ) after Kabsch alignment, demonstrating backbone RMSD $<$1.0~\AA\ across clades and conformational states.

\paragraph{S3 Fig. Per-residue RMSF heatmap.} Root-mean-square fluctuation mapped onto the reduced-system structure for each SOSIP simulation, highlighting the relay pocket's position at the 0th RMSF percentile in three of four structures (5FYJ, 6MU6, 4TVP).

\paragraph{S4 Fig. PCA variance and Schlitter entropy comparison.} (A)~Total PCA variance across five simulation systems. (B)~Schlitter conformational entropy. The G471A variant (5FYJ) shows the lowest variance (2.72~nm$^2$) and lowest entropy (1,246~J/mol$\cdot$K), consistent with steric gridlock abolishing the relay gate.

\paragraph{S5 Fig. Ensemble docking score comparison across conformational states.} BMS-626529 docking affinities: 6MEO 20-frame ensemble ($-6.81 \pm 0.75$~kcal/mol), 5U7O SOSIP 3 clusters ($-4.29 \pm 0.24$), and 6MU6 SOSIP 3 clusters ($-3.90 \pm 0.33$). The ${\sim}2.5$~kcal/mol difference quantifies CD4's geometric contribution to the binding pocket.


\paragraph{S7 Fig. Computational alanine scanning.} (A)~GNM relay--CD4 coupling loss upon single-residue alanine mutations at each of the 8 relay pocket positions. Maximum disruption: R465A ($-2.2$\%). The glycine triplet (G467--G469) shows zero coupling change, confirming the relay signal propagates through backbone topology. (B)~Contact loss per mutation.

\begingroup
\small
\setlength{\bibsep}{0.4em}

\begin{thebibliography}{99}

\bibitem[Abraham et~al., 2015]{Abraham2015}
Abraham, M.~J. et~al. (2015).
\newblock GROMACS: High performance molecular simulations through multi-level parallelism from laptops to supercomputers.
\newblock \textit{SoftwareX}, 1--2, 19--25.
\newblock \href{https://doi.org/10.1016/j.softx.2015.06.001}{doi:10.1016/j.softx.2015.06.001}

\bibitem[Blattner et~al., 2014]{Blattner2014}
Blattner, C. et~al. (2014).
\newblock Structural delineation of a quaternary, cleavage-dependent epitope at the gp41--gp120 interface on intact HIV-1 Env trimers.
\newblock \textit{Immunity}, 40, 669--680.
\newblock \href{https://doi.org/10.1016/j.immuni.2014.04.008}{doi:10.1016/j.immuni.2014.04.008}

\bibitem[Bennett, 1976]{Bennett1976}
Bennett, C.~H. (1976).
\newblock Efficient estimation of free energy differences from Monte Carlo data.
\newblock \textit{J.\ Comput.\ Phys.}, 22, 245--268.
\newblock \href{https://doi.org/10.1016/0021-9991(76)90078-4}{doi:10.1016/0021-9991(76)90078-4}

\bibitem[Capra et~al., 2009]{Capra2009}
Capra, J.~A., Laskowski, R.~A., Thornton, J.~M., Singh, M. \& Funkhouser, T.~A. (2009).
\newblock Predicting protein ligand binding sites by combining evolutionary sequence conservation and 3D structure.
\newblock \textit{PLoS Comput.\ Biol.}, 5, e1000585.
\newblock \href{https://doi.org/10.1371/journal.pcbi.1000585}{doi:10.1371/journal.pcbi.1000585}

\bibitem[Chan et~al., 1997]{Chan1997}
Chan, D.~C., Fass, D., Berger, J.~M. \& Kim, P.~S. (1997).
\newblock Core structure of gp41 from the HIV envelope glycoprotein.
\newblock \textit{Cell}, 89, 263--273.
\newblock \href{https://doi.org/10.1016/s0092-8674(00)80205-6}{doi:10.1016/s0092-8674(00)80205-6}

\bibitem[Duan et~al., 2016]{Duan2016}
Duan, L., Liu, X. \& Zhang, J.~Z.~H. (2016).
\newblock Interaction entropy: a new paradigm for highly efficient and reliable computation of protein--ligand binding free energy.
\newblock \textit{J.\ Am.\ Chem.\ Soc.}, 138, 5722--5728.
\newblock \href{https://doi.org/10.1021/jacs.6b02682}{doi:10.1021/jacs.6b02682}

\bibitem[Dorr et~al., 2005]{Dorr2005}
Dorr, P. et~al. (2005).
\newblock Maraviroc (UK-427,857), a potent, orally bioavailable, and selective small-molecule inhibitor of chemokine receptor CCR5 with broad-spectrum anti-human immunodeficiency virus type 1 activity.
\newblock \textit{Antimicrob.\ Agents Chemother.}, 49, 4721--4732.
\newblock \href{https://doi.org/10.1128/AAC.49.11.4721-4732.2005}{doi:10.1128/AAC.49.11.4721-4732.2005}

\bibitem[Eastman et~al., 2017]{Eastman2017}
Eastman, P. et~al. (2017).
\newblock OpenMM 7: Rapid development of high performance algorithms for molecular dynamics.
\newblock \textit{PLoS Comput.\ Biol.}, 13, e1005659.
\newblock \href{https://doi.org/10.1371/journal.pcbi.1005659}{doi:10.1371/journal.pcbi.1005659}

\bibitem[Eberhardt et~al., 2021]{Eberhardt2021}
Eberhardt, J., Santos-Martins, D., Tillack, A.~F. \& Forli, S. (2021).
\newblock AutoDock Vina 1.2.0: New Docking Methods, Expanded Force Field, and Python Bindings.
\newblock \textit{J.\ Chem.\ Inf.\ Model.}, 61, 3891--3898.
\newblock \href{https://doi.org/10.1021/acs.jcim.1c00203}{doi:10.1021/acs.jcim.1c00203}

\bibitem[Genheden \& Ryde, 2015]{Genheden2015}
Genheden, S. \& Ryde, U. (2015).
\newblock The MM/PBSA and MM/GBSA methods to estimate ligand-binding affinities.
\newblock \textit{Expert Opin.\ Drug Discov.}, 10, 449--461.
\newblock \href{https://doi.org/10.1517/17460441.2015.1032936}{doi:10.1517/17460441.2015.1032936}

\bibitem[Guilhaudis et~al., 2002]{Guilhaudis2002}
Guilhaudis, L., Jacobs, A. \& Caffrey, M. (2002).
\newblock Solution structure of the HIV gp120 C5 domain.
\newblock \textit{Eur.\ J.\ Biochem.}, 269, 4860--4867.
\newblock \href{https://doi.org/10.1046/j.1432-1033.2002.03187.x}{doi:10.1046/j.1432-1033.2002.03187.x}

\bibitem[Guttman et~al., 2012]{Guttman2012}
Guttman, M. et~al. (2012).
\newblock Solution structure, conformational dynamics, and CD4-induced activation in full-length, glycosylated, monomeric HIV gp120.
\newblock \textit{J.\ Virol.}, 86, 8750--8764.
\newblock \href{https://doi.org/10.1128/JVI.07224-11}{doi:10.1128/JVI.07224-11}

\bibitem[Halgren, 2009]{Halgren2009}
Halgren, T.~A. (2009).
\newblock Identifying and characterizing binding sites and assessing druggability.
\newblock \textit{J. Chem. Inf. Model.}, 49, 377--389.
\newblock \href{https://doi.org/10.1021/ci800324m}{doi:10.1021/ci800324m}

\bibitem[Haddox et~al., 2018]{Haddox2018}
Haddox, H.~K., Dingens, A.~S., Hilton, S.~K., Overbaugh, J. \& Bloom, J.~D. (2018).
\newblock Mapping mutational effects along the evolutionary landscape of HIV envelope.
\newblock \textit{eLife}, 7, e34420.
\newblock \href{https://doi.org/10.7554/eLife.34420}{doi:10.7554/eLife.34420}

\bibitem[Henderson et~al., 2020]{Henderson2020}
Henderson, R. et~al. (2020).
\newblock Disruption of the HIV-1 Envelope allosteric network blocks CD4-induced rearrangements.
\newblock \textit{Nat.\ Commun.}, 11, 520.
\newblock \href{https://doi.org/10.1038/s41467-019-14196-w}{doi:10.1038/s41467-019-14196-w}

\bibitem[Herschhorn et~al., 2016]{Herschhorn2016}
Herschhorn, A. et~al. (2016).
\newblock Release of gp120 restraints leads to an entry-competent intermediate state of the HIV-1 envelope glycoproteins.
\newblock \textit{mBio}, 7, e01598-16.
\newblock \href{https://doi.org/10.1128/mBio.01598-16}{doi:10.1128/mBio.01598-16}

\bibitem[Hou et~al., 2011]{Hou2011}
Hou, T., Wang, J., Li, Y. \& Wang, W. (2011).
\newblock Assessing the performance of the MM/PBSA and MM/GBSA methods. 1.\ The accuracy of binding free energy calculations based on molecular dynamics simulations.
\newblock \textit{J.\ Chem.\ Inf.\ Model.}, 51, 69--82.
\newblock \href{https://doi.org/10.1021/ci100275a}{doi:10.1021/ci100275a}

\bibitem[Huang et~al., 2005]{HuangCC2005}
Huang, C.~C. et~al. (2005).
\newblock Structure of a V3-containing HIV-1 gp120 core.
\newblock \textit{Science}, 310, 1025--1028.
\newblock \href{https://doi.org/10.1126/science.1118398}{doi:10.1126/science.1118398}
\newblock \href{https://doi.org/10.1021/ct100494z}{doi:10.1021/ct100494z}

\bibitem[Huang et~al., 2014]{Huang2014}
Huang, J. et~al. (2014).
\newblock Broad and potent HIV-1 neutralization by a human antibody that binds the gp41--gp120 interface.
\newblock \textit{Nature}, 515, 138--142.
\newblock \href{https://doi.org/10.1038/nature13601}{doi:10.1038/nature13601}

\bibitem[Jacobson et~al., 2009]{Jacobson2009}
Jacobson, J.~M. et~al. (2009).
\newblock Safety, pharmacokinetics, and antiretroviral activity of multiple doses of ibalizumab (formerly TNX-355), an anti-CD4 monoclonal antibody, in human immunodeficiency virus type 1-infected adults.
\newblock \textit{Antimicrob.\ Agents Chemother.}, 53, 450--457.
\newblock \href{https://doi.org/10.1128/AAC.00942-08}{doi:10.1128/AAC.00942-08}

\bibitem[Julien et~al., 2013]{Julien2013}
Julien, J.-P. et~al. (2013).
\newblock Crystal structure of a soluble cleaved HIV-1 envelope trimer.
\newblock \textit{Science}, 342, 1477--1483.
\newblock \href{https://doi.org/10.1126/science.1245625}{doi:10.1126/science.1245625}

\bibitem[Kirkwood, 1935]{Kirkwood1935}
Kirkwood, J.~G. (1935).
\newblock Statistical mechanics of fluid mixtures.
\newblock \textit{J.\ Chem.\ Phys.}, 3, 300--313.
\newblock \href{https://doi.org/10.1063/1.1749657}{doi:10.1063/1.1749657}

\bibitem[Kilby et~al., 1998]{Kilby1998}
Kilby, J.~M. et~al. (1998).
\newblock Potent suppression of HIV-1 replication in humans by T-20, a peptide inhibitor of gp41-mediated virus entry.
\newblock \textit{Nat.\ Med.}, 4, 1302--1307.
\newblock \href{https://doi.org/10.1038/3293}{doi:10.1038/3293}

\bibitem[Kong et~al., 2010]{Kong2010}
Kong, L. et~al. (2010).
\newblock Local conformational stability of HIV-1 gp120 in unliganded and CD4-bound states as defined by amide hydrogen/deuterium exchange.
\newblock \textit{J.\ Virol.}, 84, 10311--10321.
\newblock \href{https://doi.org/10.1128/JVI.00688-10}{doi:10.1128/JVI.00688-10}

\bibitem[Kriv\'ak \& Hoksza, 2018]{Krivak2018}
Kriv\'ak, R. \& Hoksza, D. (2018).
\newblock P2Rank: machine learning based tool for rapid and accurate prediction of ligand binding sites from protein structure.
\newblock \textit{J.\ Cheminform.}, 10, 39.
\newblock \href{https://doi.org/10.1186/s13321-018-0285-8}{doi:10.1186/s13321-018-0285-8}

\bibitem[Lai et~al., 2019]{Lai2019}
Lai, Y.-T. et~al. (2019).
\newblock Lattice engineering enables definition of molecular features allowing for potent small-molecule inhibition of HIV-1 entry.
\newblock \textit{Nat.\ Commun.}, 10, 47.
\newblock \href{https://doi.org/10.1038/s41467-018-07851-1}{doi:10.1038/s41467-018-07851-1}

\bibitem[Ma et~al., 2018]{Ma2018}
Ma, X. et~al. (2018).
\newblock HIV-1 Env trimer opens through an asymmetric intermediate in which individual protomers adopt distinct conformations.
\newblock \textit{eLife}, 7, e34271.
\newblock \href{https://doi.org/10.7554/eLife.34271}{doi:10.7554/eLife.34271}

\bibitem[Le~Guilloux et~al., 2009]{LeGuilloux2009}
Le~Guilloux, V., Schmidtke, P. \& Tuffery, P. (2009).
\newblock Fpocket: an open source platform for ligand pocket detection.
\newblock \textit{BMC Bioinformatics}, 10, 168.
\newblock \href{https://doi.org/10.1186/1471-2105-10-168}{doi:10.1186/1471-2105-10-168}

\bibitem[Lee et~al., 2016]{Lee2016}
Lee, J.~H., Ozorowski, G. \& Ward, A.~B. (2016).
\newblock Cryo-EM structure of a native, fully glycosylated, cleaved HIV-1 envelope trimer.
\newblock \textit{Science}, 351, 1043--1048.
\newblock \href{https://doi.org/10.1126/science.aad2450}{doi:10.1126/science.aad2450}

\bibitem[Maeda et~al., 2006]{Maeda2006}
Maeda, K. et~al. (2006).
\newblock Structural and molecular interactions of CCR5 inhibitors with CCR5.
\newblock \textit{J.\ Biol.\ Chem.}, 281, 12688--12698.
\newblock \href{https://doi.org/10.1074/jbc.M512688200}{doi:10.1074/jbc.M512688200}

\bibitem[O'Boyle et~al., 2011]{OBoyle2011}
O'Boyle, N.~M. et~al. (2011).
\newblock Open Babel: An open chemical toolbox.
\newblock \textit{J.\ Cheminform.}, 3, 33.
\newblock \href{https://doi.org/10.1186/1758-2946-3-33}{doi:10.1186/1758-2946-3-33}

\bibitem[Pancera et~al., 2010]{Pancera2010}
Pancera, M. et~al. (2010).
\newblock Structure of HIV-1 gp120 with gp41-interactive region reveals layered envelope architecture and basis of conformational mobility.
\newblock \textit{Proc.\ Natl.\ Acad.\ Sci.\ USA}, 107, 1166--1171.
\newblock \href{https://doi.org/10.1073/pnas.0911004107}{doi:10.1073/pnas.0911004107}

\bibitem[Pancera et~al., 2014]{Pancera2014}
Pancera, M. et~al. (2014).
\newblock Structure and immune recognition of trimeric pre-fusion HIV-1 Env.
\newblock \textit{Nature}, 514, 455--461.
\newblock \href{https://doi.org/10.1038/nature13808}{doi:10.1038/nature13808}

\bibitem[Pancera et~al., 2017]{Pancera2017}
Pancera, M. et~al. (2017).
\newblock Crystal structures of trimeric HIV envelope with entry inhibitors BMS-378806 and BMS-626529.
\newblock \textit{Nat.\ Chem.\ Biol.}, 13, 1115--1122.
\newblock \href{https://doi.org/10.1038/nchembio.2460}{doi:10.1038/nchembio.2460}

\bibitem[Rizzuto et~al., 1998]{Rizzuto1998}
Rizzuto, C.~D. et~al. (1998).
\newblock A conserved HIV gp120 glycoprotein structure involved in chemokine receptor binding.
\newblock \textit{Science}, 280, 1949--1953.
\newblock \href{https://doi.org/10.1126/science.280.5371.1949}{doi:10.1126/science.280.5371.1949}

\bibitem[Scharf et~al., 2014]{Scharf2014}
Scharf, L. et~al. (2014).
\newblock Antibody 8ANC195 reveals a site of broad vulnerability on the HIV-1 envelope spike.
\newblock \textit{Cell Rep.}, 7, 785--795.
\newblock \href{https://doi.org/10.1016/j.celrep.2014.04.001}{doi:10.1016/j.celrep.2014.04.001}

\bibitem[Schindler et~al., 2000]{Schindler2000}
Schindler, T. et~al. (2000).
\newblock Structural mechanism for STI-571 inhibition of Abelson tyrosine kinase.
\newblock \textit{Science}, 289, 1938--1942.
\newblock \href{https://doi.org/10.1126/science.289.5486.1938}{doi:10.1126/science.289.5486.1938}

\bibitem[Stewart-Jones et~al., 2016]{StewartJones2016}
Stewart-Jones, G.~B.~E. et~al. (2016).
\newblock Trimeric HIV-1-Env structures define glycan shields from clades A, B, and G.
\newblock \textit{Cell}, 165, 813--826.
\newblock \href{https://doi.org/10.1016/j.cell.2016.04.010}{doi:10.1016/j.cell.2016.04.010}

\bibitem[Shirts \& Chodera, 2008]{Shirts2008}
Shirts, M.~R. \& Chodera, J.~D. (2008).
\newblock Statistically optimal analysis of samples from multiple equilibrium states.
\newblock \textit{J.\ Chem.\ Phys.}, 129, 124105.
\newblock \href{https://doi.org/10.1063/1.2978177}{doi:10.1063/1.2978177}

\bibitem[Shirts et~al., 2017]{Shirts2017}
Shirts, M.~R. et~al. (2017).
\newblock Lessons learned from comparing molecular dynamics engines on the SAMPL5 dataset.
\newblock \textit{J.\ Comput.-Aided Mol.\ Des.}, 31, 147--161.
\newblock \href{https://doi.org/10.1007/s10822-016-9977-1}{doi:10.1007/s10822-016-9977-1}

\bibitem[Shaik et~al., 2019]{Shaik2019}
Shaik, M.~M. et~al. (2019).
\newblock Structural basis of coreceptor recognition by HIV-1 envelope spike.
\newblock \textit{Nature}, 565, 318--323.
\newblock \href{https://doi.org/10.1038/s41586-018-0804-9}{doi:10.1038/s41586-018-0804-9}

\bibitem[Tan et~al., 2013]{Tan2013}
Tan, Q. et~al. (2013).
\newblock Structure of the CCR5 chemokine receptor--HIV entry inhibitor maraviroc complex.
\newblock \textit{Science}, 341, 1387--1390.
\newblock \href{https://doi.org/10.1126/science.1241475}{doi:10.1126/science.1241475}

\bibitem[Trott \& Olson, 2010]{Trott2010}
Trott, O. \& Olson, A.~J. (2010).
\newblock AutoDock Vina: improving the speed and accuracy of docking with a new scoring function, efficient optimization, and multithreading.
\newblock \textit{J.\ Comput.\ Chem.}, 31, 455--461.
\newblock \href{https://doi.org/10.1002/jcc.21334}{doi:10.1002/jcc.21334}

\bibitem[Weissenhorn et~al., 1997]{Weissenhorn1997}
Weissenhorn, W., Dessen, A., Harrison, S.~C., Skehel, J.~J. \& Wiley, D.~C. (1997).
\newblock Atomic structure of the ectodomain from HIV-1 gp41.
\newblock \textit{Nature}, 387, 426--430.
\newblock \href{https://doi.org/10.1038/387426a0}{doi:10.1038/387426a0}

\bibitem[Wild et~al., 1994]{Wild1994}
Wild, C.~T., Shugars, D.~C., Greenwell, T.~K., McDanal, C.~B. \& Matthews, T.~J. (1994).
\newblock Peptides corresponding to a predictive $\alpha$-helical domain of human immunodeficiency virus type 1 gp41 are potent inhibitors of virus infection.
\newblock \textit{Proc.\ Natl.\ Acad.\ Sci.\ USA}, 91, 9770--9774.
\newblock \href{https://doi.org/10.1073/pnas.91.21.9770}{doi:10.1073/pnas.91.21.9770}

\bibitem[Xu et~al., 2026]{Xu2026}
Xu, M. et~al. (2026).
\newblock gp41-directed broadly neutralizing antibodies allosterically remodel HIV-1 Env conformational dynamics.
\newblock \textit{bioRxiv}, 2026.01.27.702099.
\newblock \href{https://doi.org/10.1101/2026.01.27.702099}{doi:10.1101/2026.01.27.702099}

\end{thebibliography}
\endgroup

\vspace{1.5em}
\begin{center}
\rule{0.4\textwidth}{0.4pt}
\end{center}

\end{document}
